%************************************************************************
% Worksheet: GPIO output stages (push-pull, open drain, open collector)
% simulated in Falstad CircuitJS. Circuits: examples/interfaces/falstad/
%
%   make worksheet      builds Worksheet_IO_Ports.pdf (students)
%                       and Worksheet_IO_Ports_Solutions.pdf (lecturer)
%************************************************************************
\documentclass[a4paper,11pt]{article}
%************************************************************************
% Common preamble for the script (docu.tex) and the slides (slides.tex).
% Each root file loads its document class, then %************************************************************************
% Common preamble for the script (docu.tex) and the slides (slides.tex).
% Each root file loads its document class, then %************************************************************************
% Common preamble for the script (docu.tex) and the slides (slides.tex).
% Each root file loads its document class, then \input{preamble},
% then hyperref with its own options, then its page geometry.
%************************************************************************
\usepackage{iftex}
\ifPDFTeX
  % fallback without the RWU font: make PDFTEX=1 (or plain pdflatex)
  \usepackage[utf8]{inputenc}
  \usepackage[T1]{fontenc}
  \usepackage{lmodern}
\else
  % RWU corporate font Barlow Semi Condensed, shipped in fonts/ (SIL OFL)
  \usepackage{fontspec}
  \defaultfontfeatures{Ligatures=TeX}
  \newcommand{\rwufontsetup}[1]{#1{BarlowSemiCondensed}[
      Path=fonts/, Extension=.otf,
      UprightFont=*-Regular, ItalicFont=*-Italic,
      BoldFont=*-SemiBold, BoldItalicFont=*-SemiBoldItalic]}
  \rwufontsetup\setmainfont
  \rwufontsetup\setsansfont
  \setmonofont{Latin Modern Mono}[Ligatures=TeXOff]   % keep -- as two hyphens in code
\fi
\usepackage{xcolor}
\usepackage{color,epsfig,graphicx,amsmath,amsfonts,latexsym,longtable,epic,eepic,fancybox,moreverb,url,version,ifthen,alltt}
\usepackage{shadow}

% RWU corporate colours (as in the RWU LaTeX style, rwudefs.sty)
\definecolor{rwuviolet}{RGB}{102,56,182}
\definecolor{rwuvioletlight}{RGB}{155,125,212}
\definecolor{rwucyan}{RGB}{0,169,206}
\definecolor{rwucyanlight}{RGB}{5,195,222}
\definecolor{rwucyan40}{RGB}{153,221,253}
\colorlet{rwucyandark}{rwucyan!75!black}   % cyan for text on white
\colorlet{rwugray}{black!55}
% pulls the first bullet up under a topic label: \item[{\rot\bf Topic}]~\\[-\lblskip]
% (em-based, so it does not depend on the x-height of the font; was -3ex with Latin Modern)
\newcommand{\lblskip}{1.1em}
% same before a framed listing, which needs room for its top rule
\newcommand{\lstskip}{0.5em}
\newcommand{\rwulogo}[1][height=1.3cm]{\includegraphics[#1]{rwu_logo_hor_lila-cyan_cmyk}}
\usepackage{chngcntr}
\usepackage{verbatim}
\usepackage{tikz}                % own drawings (block diagrams, signals)
\usepackage{adjustbox}           % scale a drawing to fit width and height
\usepackage{tabularx}            % tables with wrapping columns (type L)
\newcolumntype{L}{>{\raggedright\arraybackslash}X}
\usetikzlibrary{arrows.meta,positioning,calc,decorations.pathreplacing,shapes.symbols,shapes.geometric}
\usepackage{tikz-timing}         % timing diagrams (UART, I2C, SPI, ...)

\usepackage{listings}
\colorlet{codegreen}{black!55}
\definecolor{codegray}{rgb}{0.5,0.5,0.5}
\colorlet{codepurple}{rwucyandark}
\definecolor{backcolour}{rgb}{0.95,0.95,0.92}
\lstdefinestyle{mystyle}{
    backgroundcolor=\color{backcolour},
    commentstyle=\color{codegreen},
    keywordstyle=\color{rwuviolet}\bfseries,
    numberstyle=\tiny\color{codegray},
    stringstyle=\color{codepurple},
    basicstyle=\ttfamily\footnotesize,
    breakatwhitespace=false,
    breaklines=true,
    captionpos=b,
    keepspaces=true,
    numbers=left,
    numbersep=5pt,
    showspaces=false,
    showstringspaces=false,
    showtabs=false,
    tabsize=2
}
\lstset{style=mystyle}

\colorlet{mygreen}{black!55}
\definecolor{mygray}{rgb}{0.5,0.5,0.5}
\colorlet{mymauve}{rwucyandark}

%Customize a bit the look
\lstset{ %
backgroundcolor=\color{white}, % choose the background color; you must add \usepackage{color} or \usepackage{xcolor}
basicstyle=\footnotesize, % the size of the fonts that are used for the code
breakatwhitespace=false, % sets if automatic breaks should only happen at whitespace
breaklines=true, % sets automatic line breaking
captionpos=b, % sets the caption-position to bottom
commentstyle=\color{mygreen}, % comment style
deletekeywords={...}, % if you want to delete keywords from the given language
escapeinside={\%*}{*)}, % if you want to add LaTeX within your code
extendedchars=true, % lets you use non-ASCII characters; for 8-bits encodings only, does not work with UTF-8
frame=single, % adds a frame around the code
keepspaces=true, % keeps spaces in text, useful for keeping indentation of code (possibly needs columns=flexible)
keywordstyle=\color{rwuviolet}\bfseries, % keyword style
% language=Octave, % the language of the code
morekeywords={*,...}, % if you want to add more keywords to the set
numbers=left, % where to put the line-numbers; possible values are (none, left, right)
numbersep=5pt, % how far the line-numbers are from the code
numberstyle=\tiny\color{mygray}, % the style that is used for the line-numbers
rulecolor=\color{rwucyan40}, % if not set, the frame-color may be changed on line-breaks within not-black text (e.g. comments (green here))
showspaces=false, % show spaces everywhere adding particular underscores; it overrides 'showstringspaces'
showstringspaces=false, % underline spaces within strings only
showtabs=false, % show tabs within strings adding particular underscores
stepnumber=1, % the step between two line-numbers. If it's 1, each line will be numbered
stringstyle=\color{mymauve}, % string literal style
tabsize=2, % sets default tabsize to 2 spaces
title=\lstname % show the filename of files included with \lstinputlisting; also try caption instead of title
}
%END of listing package%

\definecolor{darkgray}{rgb}{.4,.4,.4}
\definecolor{purple}{rgb}{0.65, 0.12, 0.82}

%define Javascript language
\lstdefinelanguage{JavaScript}{
keywords={typeof, new, true, false, catch, function, return, null, catch, switch, var, if, in, while, do, else, case, break},
keywordstyle=\color{blue}\bfseries,
ndkeywords={class, export, boolean, throw, implements, import, this},
ndkeywordstyle=\color{darkgray}\bfseries,
identifierstyle=\color{black},
sensitive=false,
comment=[l]{//},
morecomment=[s]{/*}{*/},
commentstyle=\color{purple}\ttfamily,
stringstyle=\color{red}\ttfamily,
morestring=[b]',
morestring=[b]"
}

\lstset{
language=JavaScript,
extendedchars=true,
basicstyle=\footnotesize\ttfamily,
showstringspaces=false,
showspaces=false,
numbers=left,
numberstyle=\footnotesize,
numbersep=9pt,
tabsize=2,
breaklines=true,
showtabs=false,
captionpos=b,
columns=fullflexible   % natural widths, also for spaces (comments are set smaller)
}

%define Rust language (listings has no built-in Rust)
\lstdefinelanguage{Rust}{
keywords={as, break, const, continue, crate, else, enum, extern, false, fn, for, if, impl, in, let, loop, match, mod, move, mut, pub, ref, return, self, Self, static, struct, super, trait, true, type, unsafe, use, where, while, async, await, dyn},
sensitive=true,
comment=[l]{//},
morecomment=[s]{/*}{*/},
morestring=[b]"
}
,
% then hyperref with its own options, then its page geometry.
%************************************************************************
\usepackage{iftex}
\ifPDFTeX
  % fallback without the RWU font: make PDFTEX=1 (or plain pdflatex)
  \usepackage[utf8]{inputenc}
  \usepackage[T1]{fontenc}
  \usepackage{lmodern}
\else
  % RWU corporate font Barlow Semi Condensed, shipped in fonts/ (SIL OFL)
  \usepackage{fontspec}
  \defaultfontfeatures{Ligatures=TeX}
  \newcommand{\rwufontsetup}[1]{#1{BarlowSemiCondensed}[
      Path=fonts/, Extension=.otf,
      UprightFont=*-Regular, ItalicFont=*-Italic,
      BoldFont=*-SemiBold, BoldItalicFont=*-SemiBoldItalic]}
  \rwufontsetup\setmainfont
  \rwufontsetup\setsansfont
  \setmonofont{Latin Modern Mono}[Ligatures=TeXOff]   % keep -- as two hyphens in code
\fi
\usepackage{xcolor}
\usepackage{color,epsfig,graphicx,amsmath,amsfonts,latexsym,longtable,epic,eepic,fancybox,moreverb,url,version,ifthen,alltt}
\usepackage{shadow}

% RWU corporate colours (as in the RWU LaTeX style, rwudefs.sty)
\definecolor{rwuviolet}{RGB}{102,56,182}
\definecolor{rwuvioletlight}{RGB}{155,125,212}
\definecolor{rwucyan}{RGB}{0,169,206}
\definecolor{rwucyanlight}{RGB}{5,195,222}
\definecolor{rwucyan40}{RGB}{153,221,253}
\colorlet{rwucyandark}{rwucyan!75!black}   % cyan for text on white
\colorlet{rwugray}{black!55}
% pulls the first bullet up under a topic label: \item[{\rot\bf Topic}]~\\[-\lblskip]
% (em-based, so it does not depend on the x-height of the font; was -3ex with Latin Modern)
\newcommand{\lblskip}{1.1em}
% same before a framed listing, which needs room for its top rule
\newcommand{\lstskip}{0.5em}
\newcommand{\rwulogo}[1][height=1.3cm]{\includegraphics[#1]{rwu_logo_hor_lila-cyan_cmyk}}
\usepackage{chngcntr}
\usepackage{verbatim}
\usepackage{tikz}                % own drawings (block diagrams, signals)
\usepackage{adjustbox}           % scale a drawing to fit width and height
\usepackage{tabularx}            % tables with wrapping columns (type L)
\newcolumntype{L}{>{\raggedright\arraybackslash}X}
\usetikzlibrary{arrows.meta,positioning,calc,decorations.pathreplacing,shapes.symbols,shapes.geometric}
\usepackage{tikz-timing}         % timing diagrams (UART, I2C, SPI, ...)

\usepackage{listings}
\colorlet{codegreen}{black!55}
\definecolor{codegray}{rgb}{0.5,0.5,0.5}
\colorlet{codepurple}{rwucyandark}
\definecolor{backcolour}{rgb}{0.95,0.95,0.92}
\lstdefinestyle{mystyle}{
    backgroundcolor=\color{backcolour},
    commentstyle=\color{codegreen},
    keywordstyle=\color{rwuviolet}\bfseries,
    numberstyle=\tiny\color{codegray},
    stringstyle=\color{codepurple},
    basicstyle=\ttfamily\footnotesize,
    breakatwhitespace=false,
    breaklines=true,
    captionpos=b,
    keepspaces=true,
    numbers=left,
    numbersep=5pt,
    showspaces=false,
    showstringspaces=false,
    showtabs=false,
    tabsize=2
}
\lstset{style=mystyle}

\colorlet{mygreen}{black!55}
\definecolor{mygray}{rgb}{0.5,0.5,0.5}
\colorlet{mymauve}{rwucyandark}

%Customize a bit the look
\lstset{ %
backgroundcolor=\color{white}, % choose the background color; you must add \usepackage{color} or \usepackage{xcolor}
basicstyle=\footnotesize, % the size of the fonts that are used for the code
breakatwhitespace=false, % sets if automatic breaks should only happen at whitespace
breaklines=true, % sets automatic line breaking
captionpos=b, % sets the caption-position to bottom
commentstyle=\color{mygreen}, % comment style
deletekeywords={...}, % if you want to delete keywords from the given language
escapeinside={\%*}{*)}, % if you want to add LaTeX within your code
extendedchars=true, % lets you use non-ASCII characters; for 8-bits encodings only, does not work with UTF-8
frame=single, % adds a frame around the code
keepspaces=true, % keeps spaces in text, useful for keeping indentation of code (possibly needs columns=flexible)
keywordstyle=\color{rwuviolet}\bfseries, % keyword style
% language=Octave, % the language of the code
morekeywords={*,...}, % if you want to add more keywords to the set
numbers=left, % where to put the line-numbers; possible values are (none, left, right)
numbersep=5pt, % how far the line-numbers are from the code
numberstyle=\tiny\color{mygray}, % the style that is used for the line-numbers
rulecolor=\color{rwucyan40}, % if not set, the frame-color may be changed on line-breaks within not-black text (e.g. comments (green here))
showspaces=false, % show spaces everywhere adding particular underscores; it overrides 'showstringspaces'
showstringspaces=false, % underline spaces within strings only
showtabs=false, % show tabs within strings adding particular underscores
stepnumber=1, % the step between two line-numbers. If it's 1, each line will be numbered
stringstyle=\color{mymauve}, % string literal style
tabsize=2, % sets default tabsize to 2 spaces
title=\lstname % show the filename of files included with \lstinputlisting; also try caption instead of title
}
%END of listing package%

\definecolor{darkgray}{rgb}{.4,.4,.4}
\definecolor{purple}{rgb}{0.65, 0.12, 0.82}

%define Javascript language
\lstdefinelanguage{JavaScript}{
keywords={typeof, new, true, false, catch, function, return, null, catch, switch, var, if, in, while, do, else, case, break},
keywordstyle=\color{blue}\bfseries,
ndkeywords={class, export, boolean, throw, implements, import, this},
ndkeywordstyle=\color{darkgray}\bfseries,
identifierstyle=\color{black},
sensitive=false,
comment=[l]{//},
morecomment=[s]{/*}{*/},
commentstyle=\color{purple}\ttfamily,
stringstyle=\color{red}\ttfamily,
morestring=[b]',
morestring=[b]"
}

\lstset{
language=JavaScript,
extendedchars=true,
basicstyle=\footnotesize\ttfamily,
showstringspaces=false,
showspaces=false,
numbers=left,
numberstyle=\footnotesize,
numbersep=9pt,
tabsize=2,
breaklines=true,
showtabs=false,
captionpos=b,
columns=fullflexible   % natural widths, also for spaces (comments are set smaller)
}

%define Rust language (listings has no built-in Rust)
\lstdefinelanguage{Rust}{
keywords={as, break, const, continue, crate, else, enum, extern, false, fn, for, if, impl, in, let, loop, match, mod, move, mut, pub, ref, return, self, Self, static, struct, super, trait, true, type, unsafe, use, where, while, async, await, dyn},
sensitive=true,
comment=[l]{//},
morecomment=[s]{/*}{*/},
morestring=[b]"
}
,
% then hyperref with its own options, then its page geometry.
%************************************************************************
\usepackage{iftex}
\ifPDFTeX
  % fallback without the RWU font: make PDFTEX=1 (or plain pdflatex)
  \usepackage[utf8]{inputenc}
  \usepackage[T1]{fontenc}
  \usepackage{lmodern}
\else
  % RWU corporate font Barlow Semi Condensed, shipped in fonts/ (SIL OFL)
  \usepackage{fontspec}
  \defaultfontfeatures{Ligatures=TeX}
  \newcommand{\rwufontsetup}[1]{#1{BarlowSemiCondensed}[
      Path=fonts/, Extension=.otf,
      UprightFont=*-Regular, ItalicFont=*-Italic,
      BoldFont=*-SemiBold, BoldItalicFont=*-SemiBoldItalic]}
  \rwufontsetup\setmainfont
  \rwufontsetup\setsansfont
  \setmonofont{Latin Modern Mono}[Ligatures=TeXOff]   % keep -- as two hyphens in code
\fi
\usepackage{xcolor}
\usepackage{color,epsfig,graphicx,amsmath,amsfonts,latexsym,longtable,epic,eepic,fancybox,moreverb,url,version,ifthen,alltt}
\usepackage{shadow}

% RWU corporate colours (as in the RWU LaTeX style, rwudefs.sty)
\definecolor{rwuviolet}{RGB}{102,56,182}
\definecolor{rwuvioletlight}{RGB}{155,125,212}
\definecolor{rwucyan}{RGB}{0,169,206}
\definecolor{rwucyanlight}{RGB}{5,195,222}
\definecolor{rwucyan40}{RGB}{153,221,253}
\colorlet{rwucyandark}{rwucyan!75!black}   % cyan for text on white
\colorlet{rwugray}{black!55}
% pulls the first bullet up under a topic label: \item[{\rot\bf Topic}]~\\[-\lblskip]
% (em-based, so it does not depend on the x-height of the font; was -3ex with Latin Modern)
\newcommand{\lblskip}{1.1em}
% same before a framed listing, which needs room for its top rule
\newcommand{\lstskip}{0.5em}
\newcommand{\rwulogo}[1][height=1.3cm]{\includegraphics[#1]{rwu_logo_hor_lila-cyan_cmyk}}
\usepackage{chngcntr}
\usepackage{verbatim}
\usepackage{tikz}                % own drawings (block diagrams, signals)
\usepackage{adjustbox}           % scale a drawing to fit width and height
\usepackage{tabularx}            % tables with wrapping columns (type L)
\newcolumntype{L}{>{\raggedright\arraybackslash}X}
\usetikzlibrary{arrows.meta,positioning,calc,decorations.pathreplacing,shapes.symbols,shapes.geometric}
\usepackage{tikz-timing}         % timing diagrams (UART, I2C, SPI, ...)

\usepackage{listings}
\colorlet{codegreen}{black!55}
\definecolor{codegray}{rgb}{0.5,0.5,0.5}
\colorlet{codepurple}{rwucyandark}
\definecolor{backcolour}{rgb}{0.95,0.95,0.92}
\lstdefinestyle{mystyle}{
    backgroundcolor=\color{backcolour},
    commentstyle=\color{codegreen},
    keywordstyle=\color{rwuviolet}\bfseries,
    numberstyle=\tiny\color{codegray},
    stringstyle=\color{codepurple},
    basicstyle=\ttfamily\footnotesize,
    breakatwhitespace=false,
    breaklines=true,
    captionpos=b,
    keepspaces=true,
    numbers=left,
    numbersep=5pt,
    showspaces=false,
    showstringspaces=false,
    showtabs=false,
    tabsize=2
}
\lstset{style=mystyle}

\colorlet{mygreen}{black!55}
\definecolor{mygray}{rgb}{0.5,0.5,0.5}
\colorlet{mymauve}{rwucyandark}

%Customize a bit the look
\lstset{ %
backgroundcolor=\color{white}, % choose the background color; you must add \usepackage{color} or \usepackage{xcolor}
basicstyle=\footnotesize, % the size of the fonts that are used for the code
breakatwhitespace=false, % sets if automatic breaks should only happen at whitespace
breaklines=true, % sets automatic line breaking
captionpos=b, % sets the caption-position to bottom
commentstyle=\color{mygreen}, % comment style
deletekeywords={...}, % if you want to delete keywords from the given language
escapeinside={\%*}{*)}, % if you want to add LaTeX within your code
extendedchars=true, % lets you use non-ASCII characters; for 8-bits encodings only, does not work with UTF-8
frame=single, % adds a frame around the code
keepspaces=true, % keeps spaces in text, useful for keeping indentation of code (possibly needs columns=flexible)
keywordstyle=\color{rwuviolet}\bfseries, % keyword style
% language=Octave, % the language of the code
morekeywords={*,...}, % if you want to add more keywords to the set
numbers=left, % where to put the line-numbers; possible values are (none, left, right)
numbersep=5pt, % how far the line-numbers are from the code
numberstyle=\tiny\color{mygray}, % the style that is used for the line-numbers
rulecolor=\color{rwucyan40}, % if not set, the frame-color may be changed on line-breaks within not-black text (e.g. comments (green here))
showspaces=false, % show spaces everywhere adding particular underscores; it overrides 'showstringspaces'
showstringspaces=false, % underline spaces within strings only
showtabs=false, % show tabs within strings adding particular underscores
stepnumber=1, % the step between two line-numbers. If it's 1, each line will be numbered
stringstyle=\color{mymauve}, % string literal style
tabsize=2, % sets default tabsize to 2 spaces
title=\lstname % show the filename of files included with \lstinputlisting; also try caption instead of title
}
%END of listing package%

\definecolor{darkgray}{rgb}{.4,.4,.4}
\definecolor{purple}{rgb}{0.65, 0.12, 0.82}

%define Javascript language
\lstdefinelanguage{JavaScript}{
keywords={typeof, new, true, false, catch, function, return, null, catch, switch, var, if, in, while, do, else, case, break},
keywordstyle=\color{blue}\bfseries,
ndkeywords={class, export, boolean, throw, implements, import, this},
ndkeywordstyle=\color{darkgray}\bfseries,
identifierstyle=\color{black},
sensitive=false,
comment=[l]{//},
morecomment=[s]{/*}{*/},
commentstyle=\color{purple}\ttfamily,
stringstyle=\color{red}\ttfamily,
morestring=[b]',
morestring=[b]"
}

\lstset{
language=JavaScript,
extendedchars=true,
basicstyle=\footnotesize\ttfamily,
showstringspaces=false,
showspaces=false,
numbers=left,
numberstyle=\footnotesize,
numbersep=9pt,
tabsize=2,
breaklines=true,
showtabs=false,
captionpos=b,
columns=fullflexible   % natural widths, also for spaces (comments are set smaller)
}

%define Rust language (listings has no built-in Rust)
\lstdefinelanguage{Rust}{
keywords={as, break, const, continue, crate, else, enum, extern, false, fn, for, if, impl, in, let, loop, match, mod, move, mut, pub, ref, return, self, Self, static, struct, super, trait, true, type, unsafe, use, where, while, async, await, dyn},
sensitive=true,
comment=[l]{//},
morecomment=[s]{/*}{*/},
morestring=[b]"
}

\usepackage[colorlinks=true,pdfpagemode=UseNone]{hyperref}
\hypersetup{linkcolor=rwuviolet,citecolor=rwucyandark,urlcolor=rwucyandark}
\usepackage[a4paper,top=22mm,bottom=25mm,left=22mm,right=22mm]{geometry}
\usepackage{titlesec}
\titleformat*{\section}{\sffamily\Large\bfseries\color{rwuviolet}}
\titleformat*{\subsection}{\sffamily\large\bfseries\color{rwucyandark}}
\usepackage{fancyhdr}
\pagestyle{fancy}
\fancyhf{}
\fancyhead[L]{\color{rwugray}\small Embedded Computing -- Worksheet: GPIO output stages\ifdefined\solutions\ (solutions)\fi}
\fancyhead[R]{\color{rwugray}\small RWU}
\fancyfoot[C]{\color{rwugray}\small\thepage}
\renewcommand{\headrulewidth}{0pt}
\usepackage{booktabs,enumitem,siunitx,amssymb,needspace}
\sisetup{per-mode=symbol}
\usepackage[most]{tcolorbox}
\setlength{\parindent}{0pt}
\setlength{\parskip}{0.5ex}
\setlist{itemsep=0.3ex,topsep=0.5ex}

% generated by examples/interfaces/falstad/make_circuits.py
\def\linkA{https://www.falstad.com/circuit/circuitjs.html?ctz=DwYwlgTgBAZgvAIgIwKgFwM6IEwFMC0ADAOypgiJKEB02qIARogKyGoAOTCrqAbhC1QBbTJQIkEAPgBQUKMCFCoAOyGIA5uzAB7fKu0ZU8BGyi80iGgDZUDXJeqEkdKBi6mAJhgCel2x8pbCxNUNwCQgHoZOQUlVQ0tXXYhAyM-M2DrW3sTR0IADlD3VC9fEKgGcJQKzKLwtijZeQAPKGbEABZ8qABmbChCAbTyjDAcUzQcgEEALigABQBXDAALfHZFgBtNqABxeYBJAHkobUW0DbQoEBWAQwhbkEmIMAw0cigAChBFiF5cdAPEC4CAASgQjRirXaCC6UCsHQGQ1g6VGlA6qEmiAAMtp1B8wMpLlA4KdzsSILh8W8QRUwGhqFAACorAEYADuuFw7FcZwgwNg2n5uAwUAAap9DgA5UGwCDaJREajMcXobS9ajdMXUCHRFptTrdfLdQamYymNHIDHoHLY3AwK4YEDadi4OYADRJ4slBxlABooABNL1aZTXX6U5RXT4AMnUaAA3AM5mcrtoYOhWVBQ6DGQAlMDqFaO52uubsbSc6CEhYAWSOAGUoLdlB4oFL6w2dZD5NiDchsN1sMwrFB8oNhzYUeV2IEoCsxggetQeqhNto-D3gAd+84hyOoEhiKPJ8MLZsHMxUAvEMvV1uMzDsB0JwVeoREdgCsNqikEjp1hSQwH37Z9Bh6D930RCDrXNYQNwQTQAP0YC9WAdlQJfKAvyHLDT2nBo0Iwp88IPMCoLPXUmnQ-sYOw58oOwkdKK3Yjb0gydGMHQoCKomJ1FoyC+nAoTmBcOCtzzQTESQABOfo6KQKwzXSdk4KgIRbmaXhbxXPj5DY5BjyYk9SKnCSiOkkzeisfdzM3NDNjDGEels6zXKHZ9hmtLEEGlfTgCmRRcGeWi3M4joERM4ZVw0nILQoSI0N4ftIug-JETS6LeKgNTbyMAQEGHUxNO0nBqCQK86VuLBkGoKwry3ASYSynoMqgLKOiUli0PVXBlEoUhXFnZAf1k60RrYPUYh5HSRkSjEexmqAcm6sx8tcRKaEIRq0IibRepWgbuFCEaeOMFwRpQab5FmxAhqdBwUCW26VsQKc5oeraKoC-aZGACJwAgGQgA}
\def\linkB{https://www.falstad.com/circuit/circuitjs.html?ctz=DwYwlgTgBAZgvAIgIwKgFwM6KQUwLRIAMqYI2hAdAEyogBGiArMVAA4MLOoBuETqAW0xN8SFAD4AUFCjABAqADsBiAOaswAezzLNGVPAQtuaRJQBsqOjjMVCSGlAwcWAEwwBPM1dfYrpo1RnX0CAeikZOQVlNQ1tVgE9A28oE1tLKGtbQkIADiCXVHcvQMyQlEyAlmDvcOlZAA8oBsQAFlyoAGYqKEJe5NKMMEQqFjQbBAAhAC4oAAUAVwwACzxWBYAbDahNBbR1tCcAdzA0EGWwRVUoAEMoEBvWG-A0MG4cKA3NG9cEOsimi0EO0oOZWr1+rAUkNsK1UONEAAVZYfDAARwWNwgHyON3eUDAGHQKLYEE0qggNwUaHJqg2l2uaBJGkUAG5iR9WNjXBA3jhoBgNjgjlBKeMnDhMFBlpoRTAbhhDkyPqobuMiRgTmdlhQ-hFGs02h1ch0+ixDNVhsg4egJgAZHAwQ4YECaVg4Waq8UAcVuilc8wAkgA5ChQABKYFUy2drvdswwC1YrA2HnuCwg2MUhxu-oJUrx-JuqhwYYAyvTXPyiZpFBzRVGY7r-rJw4aEABOcxdUZQEHdc0pI6GRwwPjIHKTygsAQ3BrcbAUcyMKxgBWL5d6+rAEVArs9voODoDgbEFvAQPto8HqBIADs3ZPUMGGxGqAuiE6FE6W8iu+wD43qMx6jKev6yBsdZAsBQF5FAVAmgMNoIgg3rgcAMDtlQrR9FQd49J0hDgnho5+FAiSxFoayJPo56YdBOFdOYx5EUx+TPjOmiUdoui0fqO5YYxJHwYxT4Wuh-4INhfQDiJMnMWB56SYRxEml0rFibU-HKRpvYqV0jCkWE-HXEC+mdLk4L6a0SAZOJSlYXBsnSQein8W2Zmsfej5ea0NriVAw4pLO86ft+6EAILyJK-Ltvp3nqcRjB2Z+ggTNUZDGduOlWXpzGuRxElxflsmtN0BX2fxkHFSB8kdAh7GGMhExzCG6EgO2ZUEb2XV9jkYH3NgohniZnXlThfS9V1bnbj8ABWUATFQFQ4FaLA4IoiAAMKPM8py5iAHwABQAGIAJSCJtyCiI4AhArgeCEB2QQBHa3wBg8TwvAdHxbQA+naqBfKofjnvNi2pYta2oBtiAndiGIbSAaZHQAEgAXhd5FXUQU52IIQLMFOOT404ASIrSQqwAjCxIx4QPkqD-HgxMK5QyksMIBWwoRmqx0AGqhIoGBYwIV3ThUd38GTiBzNyvLvAKQoimKAtCyLDMg8g6GaItV13kErCfkZjhG4E+qRKwqTQplcItpbENScYNttOhoSaPxuuc-ehvGwMSCdEQd4-mwfgW7IVsLggIcutkK7nu7UjAKE4AQFIQA}
\def\linkC{https://www.falstad.com/circuit/circuitjs.html?ctz=DwYwlgTgBAZgvAIgIwKgFwM6KQUwLQAMAnKmCNgQHQBMqIARogKwGoAOjCLqAbhM6gC2mbPmIIAfACgoUYIMFQAdoMQBzNmAD2eFVoyp4CVlB5pEVAGyp6OC5QJJaUDJxMATDAE8LN99htzY1RXf2CAemlZeUUVdU0dNkF9Q19TIKsbO2MHAgAOELdUTx9gqHowlHKMwrDWSJk5AA8oJsQAFjyoAGZqKAJ+1LKMMERqEzRsgGEALigAFQB3LSg2AFcMAAs8dYAbXagtNbR1zEOlc5woXbAlK4AKeg2oEC0lSfftJQBKBAbolptBCdKCWdr9QawNIjbDtVCTRAAIS0aE2q1uGCgAEMIFdFpAcO50Fo1DhUTgIAAaKA4ShqSjoZZQdw4HhkHCYt7YlybHGE663K5aaBY9EXLbCyZEtArUV3MBqTb0I4QSh-KLNVodLp5LoDExGEww5Bw9DZADKZKgAEEoHAoAAJbFKImIu1QAAyUHuwtM7NMFIwWO+DPNrzYODmNzuLzWEFx72pbC0iwpUFu6E2VwACgBZADy5sOMBtzulWagADkC0WtCXEWr-nIvUCnF1qExLFA8gMO9YoWU2AEoJtRghupRuqhdlpfE3gABJLXIajtztQJAAdi7faGRt29iYqFHiAnU-nJaB1HavfyPQI4PGBQHVWS8W0O2SBgvy+vA26D73uCAGmoaQizggGgfno34asAiy-jeUBPshSG7gO9RwQhV5oeuf5AXu6qNPBy4gch15AchnaEfO2GnoBfaUauz5gfOaikYBvT-pxTDOKxcEAEoceCSBEH0ZFIJYBppIsYFQIIWJNDwp6TkR0R0cg25UTuuH9vxxEtswq7aVAnS9tRGHsGkJ7jqp1wQZhxFLkCvFrl2TBaehckYAeORHiOY5nmpciXogYK3l0m4MfkQyvhBUGJF+wXAKFCDhZRUXAQ+hHyfFCS6EltHLulKElRZ+nqcVumgkhIE0Vhy6ZeR4JNV5c4NUCrXrq1ur1cR7GdZx4xQE13S8X10RCYNIliSNgGSdJZSyWkClKSp54dcwnnrmVentcRGlkYx3TdDu5X7dEuwXECJ1nV2t3NUMpoIgg2YLpWAD61rJYdp0mQ97SSRNcjWgoZJpjdf2A+5xnQ0MU7ydkRrkBEm1cLDkkjVDQOWUVnXY12m5-W1qMHY1xPdQx52k5d12IETd1Yzu15PfC2RvZ9iLJSsOBKGMz4YEOyB7cgVnBBq0RsKYp4hCjVAEEe87hFocE83zXAhELONGM4WuSI0kvSwgm6y-YKBNob2Rbrw9Omzk5twcr0jAOE4AQNIQA}
\def\linkD{https://www.falstad.com/circuit/circuitjs.html?ctz=DwYwlgTgBAZgvAIgIwKgFwM6IEwFMC0ADAJypgiJKEB02qIARogKyGoAOTCrqAbhC1QBbTJQIkEAPgBQUKMCFCoAOyGIA5uzAB7fKu0ZU8BGyi80iGgDZUDXJeqEkdKBi6mAJhgCel2x8pbCxNUNwCQgHoZOQUlVQ0tXXYhAyM-M2DrW3sTR0IADlD3VC9fEKgGcJQKzKLwtijZeQAPKGbEABZ8qABmbChCAbTyjDAcUzQcgBEALigAeXZcZXwPCABDMGUKgFcMKAAKAElsAGF8ABswAGtcAEozXAgMPah2PYALfHeLi4RGmKtdoILpQKwdAZDWDpUaUDqoSaIKa4XjkXBQJCuZZeKDrKAgC7aEDXAA0UA8KLRUH6IHW2w+2guHnQH3RDFehIA7odOZBcMyAIIAOSmd2oUAACjtfvgduwoAAlAD6EtxymZ7P2tPY63AaDpIHRpyVmtxEHRGCuFOe1H+0RabU63Xy3UGpmMplhyHh6ByABVWW9Pt9pRdXPr1GzcFzyRAwLxcPs0IGMOshOjtbqwPrlIbYNpoCBtEIdXGMNplOKAMpFpYYuYAIQAqlW1cyJRLq7X0dg5iAdhBzco0FAtoqVbaAfIFY6EMQrNSXVBQdgXcNTJzjC4YAJkIR9-vHMJ1s1eJRqFZmLYwOssMgL1ep8BucD54vus5uqvCtDIvbgEcs6fu+GIAOwLt+66hBcDhXlAHxjAgPTUD0dpNM+QHgSB2AdIMkG-g0-4wLOOF4UuPSEBC+EesI2gJDoegpIYT7qLOFEQuxvSUb0PQ2ARaGAiROG9FYrpQjRriIUgPqIggFKonmKBPgAMrOuFkd0PTMP01HpOwgTwYhyGoVAhJ+E+gHAupIEdCua78a4MG5HBCGIMZAnyMRVl2d0zDOO+UFQCk9G6PozH-qxwJ+f0nHRTxfE0U+QKdH0Ilie6MJSTJOTyVSdBPjOVn7hitnLsVSDOIFm7pEIJ5nkhKEecA0BFYM0kQtZ2DMAl6S7h0oH7k1L5udx2CgTFo3dVBT7DUhk0Lp1U38TNanFV1C5xet03-rNm3dVAzBYVty07bOh0Qft4J4UtiWnVF-ljf0e09X+6EXNsUVHft53UjhwzZYgzZVk1ICzld1KXbh8WBRQCB4PgVBNZFiDgz0vFglDvE-rd6GzW+35QG+0Xbbjs5E-5wHEyd6GWZQq4HRTWFUxJlqwagrkNahT5eTgUMdOCXEdeCwzVMFCCaAxySpNzJFQ91gyxVYGXlGLEuhUxQ2y4M-NUXLlXUzEs2kQzOly0rJOG2x3FdAr3HM+Zd0jRCcWxbZFvyMjc3O1hnFK9jDvoYVTvLqCnEdEgL0bhJtWnm5jUrcCSBMw9evbgHlvArF-k-fbr0Zywyf9ODueEW9H0Fxt2dYTb-0IjkHYg2DqcY4MzDEGn5Sw-DiMsU3rfty3YJ-QbDrAtgwng26gVetJdeIO8GBfD8YbmjATzLIaTXrB4ABWUA5MPuCIaYyyIKc6w6nqBrogcABidzCMoYgI6LifiKQ4aA68mZX7mRomnsVAhJ1CBCfNvPeORw6oCPukU+CBz6X2zNfQ499H7PwqsIN+RAP4YGCEA7QIDkChDvKXGI4D96UD4jA8ocCFSJjALg5BBxtAfCEBgB+QUn5w2wKYIQVkBqDU-ggKUMo5TjglPgwhSl-zkJyCZahJ8uG33NAARx2BvbwhwAASAAvDhQguFUAPEeIKY8DzGM9MEU4hJiSwFUeo3M3hJGgP-NofeXDsAmQwPpZAL0XA+LYPaGI8p6qelhvCKcwSKEglMKE0I4SmoRG0K49xlDQg+IjiLDggQgnyBCW5eJDgnCJOSdIYAERwAQBkEAA}
\def\linkE{https://www.falstad.com/circuit/circuitjs.html?ctz=DwYwlgTgBAZgvAIgIwKgFwM6KQUwLQAMA7KmCNgQHQBMqIARogKwGoAOjCLqAbhM6gC2mbPmIIAfACgoUYIMFQAdoMQBzNmAD2eFVoyp4CVlB5pEVAGyp6OC5QJJaUDJxMATDAE8LN99htzY1RXf2CAemlZeUUVdU0dNkF9Q19TIKsbO2MHAgAOELdUTx9gqHowlHKMwrDWSJk5d1jVBDB3HABDABs8CBw1VLKwLGR8JAAWVAhRkyV7AiYbHl8G6IAPKHXECbyoAGZqKAJjoZMMMERqEzRsgFEALigAeTYcJSh3CE6wD4B3MBoAAWUE6UCYUAAalA2ABXbq9WFsJ6CTrrMCCWGKTQfHhabpoTpqHAINZyTbbBC7KCWCbHU6wNIXbBTdDZABCWmBUAAcgBZZ4AZVB-RcALQICBOHc5S86ClLk6ghwLjAaiUPUoUAACgBJHlIKAjUHKLQQVHdA6UfZQqAAcT1zyNHy0byUeC+Pw+yQ6ABpSVFyVsdns8nsTiYjOdLshWbdEHqeUcwRDIXg0PicN8lGgoK73h7vr8Yb8tQAVBX9NTwzoZ6AgToupTdOUYLSwiAgFUgDv9HNQAAU7m0HQAlJQA41gDxgwgAJyWA5MI4LqBIOfOKOoP5bqCo9YrBD7a2T6JqWer-bLqCXvKsrdk4Ayymr9dHJBERdvs5CLSIdpdL0VankGlLUhMSCLhGP6qiyqDxke1q2lWNZ1iBwB-LOH5fhuUDXHs36MhEgYYbO+Frrh1BEO+G4-o+mGUlRNFHPs+w4ZuqwkQxiCsexUATAQJyEQ+JHdB8lLYRR76fmuliRjs8HZJCAAiynoSAZH5FJeFaUxMHkGMhDWI+56Mbp1E6Xs+wEPenFTtAjEWW+eEWXpREmPwR7WYJ6GmVcrkBSxNl0SRFJXNQdK8fS8llMysaKYgfIAMIAKpQN0WidO46EAIIKDgtwOTsgl4Uwi4CUJtFEfsQjZOcBn1CRMCzlFTH8SVblGFUyTxNoeBJCk9EtWxLksSNwl2dEzWUlFbF7BVByWAU7m-r1Oh6AYJmzgt+x3u1JwQdYK1DWBHVhvtBzXCFU7cVSJWHAd93Ltd0S3Tt1z8bSl0xY1U5idtX0PZ9kXPURcbZImKCPgASgDdIRQdX34TBO5pPuh5LI+RVUkjAnA99MGeRMRA+TDZF47slVHJTKNGJunlIIJgkOEIaKHkglCWEs5RgJ0owc1z6G6uTB15Cc1DnTTK0uN09jc0CMbHjVJ1XJLYsHCNUsiTdw3lers1LS9ci3VFlPgrhWuTcbs5MBb6u2yxhvHU1Nu4XNNJlYty27j1CAaH1G2+bOXOLrtdIh-xkFG6RlIR1eRxx1dzs67HntA9RJwPdH-2Uhn31QHn8dDODCb6rQZOp4uCMe1X+S02jbMCFjwee9Xic-dMOwk79GwtTZ4ITHS0HS3FkwJQgShmha9qOr6eZuoWXp7loHToWFXAyUwg-RTBo8l1wULppm2a5vm7qesWOLoVlABWUDZFRqA4DGJjvIg0M4BcGCEkoXaDloQJBAYFHEIeYVISZCDAt3c4QRtTwkRGwDA6AtDgihKgDKagAiPlvvfYqT8X5PzAR-L+P8-4DgAUAkBe4wHExMIIKBTMQhBHQVoTByAQizGvu4O+2Qx73wIffIhn8RikJVOQwBwDQEUB8nuCSTMmYsxcEEZKaUMpZUHErW04QoC6gAPqqSoRgrBJEUFvwQNQbmGA2CIFpEMZw1jgiBmiGwUwTIDKYycXIFx2Rdq8DcU3Ei4QtAmMEdgKoVibG2XMewAInjgAuMPPVAJU4gkhLMR+EIDjbFESqA42gcSEk8RCA1XIUNAnBKkMAcI4AIDSCAA}
\def\linkF{https://www.falstad.com/circuit/circuitjs.html?ctz=DwYwlgTgBAZgvAIgIwKgFwM6IEwFMC0ADAOypgiJKEB02qIARogKyGoAOTCrqAbhC1QBbTJQIkEAPgBQUKMAAmQqADshiAF64VuCPmwAWAGqp4CNlDBZk1YkgMBmZtgAcSAJyvCLnw-GkoCGsLFRxUBl4cAwQAehk5YAAPKETEAxcoB2woQhzTRAsMMBwLNFxEADEALigAeXZtfBAAewAbVtwQNGboBQgwXl1YHqgAQygkbKNA3FbRgE8oFrBWqBG5oXZY+Plk1IR0qAA2Axy82ALUIspo9HKEABUAC1xM6gcoaYBxAAUASVqUAwAHcwGgQC8MGMVFAAHI-WHoCCjFRFDDdaAACgYuDQ4zgE0IhAAlOgnhBmgBXADmTygACUAPoAIWoUGer1azVGCiBs06GKWT2azQwuChEFmCyWzRWa2gG3Y1G2sl2KTSGR8ZwsZkKxWQtzKiAAwlyxVAAOQwVrzBijEAAawt6GaYwUvLQL1gkoIwJeszAKmpUAUYGaClwAG41g0YRaGPN2KMMBhnd0oOxKWgxlBDJ8oFodNBA3z+uKoKDPZY0Mq4qrgPT1QciRNXFADC33Ed8uZUMDdcJRolIsg6HWEsCmx3cpMMgYTq2XD22OP5JP9vPTrPjnnZ8uVQkAIJCIS4oYblvb6cTYjdi4IBzCe6FCi91fAVYXme39stww6sIoFaApqBQd9oC-XMO1-XIHEIW4B0CNIiQPNcp0vH8HAcI4bzvAd33XRAsJwuDTmvUj93fEAm2I3DMmwzJ4J7OgljEIh3FQpIm2wbwoNyfj9yBfV7FQI0EElOZFmWVpOL2EoMmwBcBPvPUblE59QKgAAJAAaGC6ieIROMI5tYPg9sFwolTON5DdLPMzc+MEoRmkoWEDHsJd3w0KcFyQH9HP-ZzXIQQtdH0YxjN805-ygZwFI7SidmAEz4roxz-Lwy53yhfY0titLMsE9hLigFzsuSuTuAMGKkBw5TEOuA11MQBMkxTTjcrSPyAr81xisoYQQpXSqpxcGcjgyBrSqakS7kQa1bXtB1OIgxAjl3NsNtOKzEIEBBsA498qpOXJsDq7VBNmw17kVVt8xiOKoAAdU4gBlJtmEmfSvuyYiBt7MrhtQVp9VKJ4wgIqcu3onCDBhrIALfZKTPhkieP0xGkvrVGW0RzGmOsqG8rbWi0qxomUc+0mMe2zIeOxicmzp-GWcJ-DkuzDdDhqjJfvbdJitCZBwnuKgLHKhAIxgUZKVaNAory77HGyfmPKy5Gcc+5WLrVlDKfrVoYSV7J1bigK6p7G6TU44MTYs041a7Rn5EbfZsGYOHDh4udBesiszBYmB9p4iWhxHJBqCOZhwjAZNKCjmPwO4z2Bd99PBP28WRvrHkACsoHuBwl0LsHUG0RB6XFKw8RUEBXkxZpDIwYlhGF8XhH2bOiWoQo0EQZlk1eSU0UFJlmRB5pqUG9988LoiAlwMvC+FipJQAR0pbQQEWTEtI0VuyvbiWu6JPvEDeysIUDYNg9wTft-mSfp5F99XQrh8Y6BEqEBOZiOEuDsBI7AoAjhfCwHOCQYjNGSu-YWd4MA-z-veFAGZBpAPkCAkcj4gSvhoIQJOyVoGwJXkRFiiD1oIVKj-OgGDgBYMQAEDAr4wJEJgdIYAMRwAQBkEAA}
\def\linkG{https://www.falstad.com/circuit/circuitjs.html?ctz=DwYwlgTgBAZgvAIgIwKgFwM6KQUwLRIpRgjYAMAdAEyogBGiArGagA4MLOoBuETqAW0zZ8hBAD4AUFCjABAqADsBiAOaswAezzLNGPDACGGNKngIWUbmkSUAbKjo5bFMgGZUGDpYAmGAJ62jj7YjjYWnnQhEQD0UjJyCspqGtqsAnoGxqawQVbh9o7OFq5UjJF5foERUFGhtQWR0Sxx0rIAHlDtiAAsABxQblRQZCNmeRhgiFSWaMUA4gBcUABi2VCsAK4YABZ4WwA2B1A4Pqo4GFCailCGUMxQAPI7CgAKAMIAQugQhiA4CFaCU63QQ-Sgdh6IzGuRqk2wPVQc0QABUdjgfn8MQAKJBKS4+HAHQz+AA0tzomk2aCgSEYUBACmuqwASj0AJQnRR+W40tDo4iKLY0zQwW6KTT8nDQEA7MCsZb01grDYHbZQACiAGUACJQHxaQkYCiA+IdLq9AZ9AajSzmSzw5CI9DFR7U4WC4iYWDZC40jA4NBoMCKVRQbEsgD6zMMlOptKoBHpzwE7IoUC1BzAhIgGGWAYgYAuUAgFzAJk00CjWvJGAOOAA7iXDHN01qQJpWDhlq8AJIAOXFPigLIAGiagbIWRaEABOOxQMoL8FL8Y1BvmGiwPjIMh7-eUSwCQztbjYCh2cq1MDGc+X01tYBN0HzxeMBdIKgDVewlpm4C9jOn7fu+tIAOwLj+9qeAc-BQHKiBuBQHiTsAYqglQPSjDMAxuGQUI4WuRAZCkWj7Bk+hGCYD4JOh0xYYM+GMVCeHOtBUAkQg6hkbolHZDRsjPvR2FkN+DFQUEqFCQgmHYaBsnMWuf6PtJrGLphzFvg4v4CU+M5qUumlfn0Sm6WGoJqUMoyWYwW7Qah04WUxSCzsMalIHYdp5BueTHqeiHIbp0lIBBWmLuJ76mVJ+lMYZbhuJBkU6ahBw3BZCVhfFkGYWuzrIggfb9rp0DpYly5DFppkli4KCoWgxygj0FWGXY1qVTp+rFLgeBkLOqBgConCWA2UwIEMZkzk1wzGVAU2DDMUX-uZiCtd+bWrfNXmxP+0kbYZYG2e19k7TOB3TaBfQeUdkknaCl1lVA1pydpx2Pqlp2Hftn05bCeXFGOukgDOT1hSDWVVaQnCiDQqHLQgYMZWDrWLY+w53Vdhn3bSrlVZxhJGJsBymA5wNXS5wxYx5W3DexflnmNgWoWjiBY1Zj0Y0ltOaIg+OGITxNLaTC5s6ztko8Cp0ZVhtowuxjpIH9iClv8YDcNKgrCrphg+AAVic2BgagOCjZYOCKIgLJliYhiKP84aaC8GDsoI5slMRoJIHunjhFq0pFpcpaTBWVaRlqWu6-rnBGybRuu5mjYji2OIAGoxIoTsuzVgigleJiIG6aAenWCe-HM4ap+nztQAcmiqKEqGaFy2DaRgrArWuW5txEZoJKwVgTJD5STr3kc-fTDqD7pMSaP+jdm9MhtQK37ewkQXcoD3sh9-THhL5Dh5D5vwB98UVB9f3Y2ePvriH4+09SMAMTgBAUhAA}
\def\linkH{https://www.falstad.com/circuit/circuitjs.html?ctz=DwYwlgTgBAZgvAIgIwKgFwM6IEwFMC0ADAJypgiJKEB02qIARogKyGoAOTCrqAbhC1QBbTJQIkEAPgBQUKMCFCoAOyGIA5uzAB7fKu0ZU8BGyi80iGgDZUDXJeqEkdKBi6mAJhgCel2x8pbCxNUNwCQgHoZOQUlVQ0tXXYhAyM-M2DrW3sTR0IADlD3VC9fEKgGcJQKzKLwtijZeUUVNQQGDAwkAGZC2HTzSmpu7Idqt3TS9MrAmvSwv0aYgA8oZcQAFnyobuwoQn208owwHFM0HIAJAC4oABFcE-VldABDDABrKAAKV4BKW7MKAANSgAElsABhVy4ZQYbTQOwACzAyg8UFeUAAQgBlHE9bYAG1wvFwhNcKJgFwgCCW8lW6wQWygVg2+0O-WOp2QG1QF0QD145FwUCQP26wxBUAAspCAKoAGig2nYsKgHggr1RfxhaIw6CRIo2hAOH0uAC8oCBCdoQB8ANwGkUYWHw6DfIHApUqtUarXKHUgV4vJHaQnotCGqCE1Eim0Ad2otOi9LWm22+W2B1MxlMJ0ovPQOQASgB9AAKGLRUEhpYYAFd9doXpGRZ7XGAPCLvkGGMSda8IM6Y12IBhqFAcSAfaLbji7gBBUvMUHB9HzpfdYESyfT1VQbC3ED1iBD5RoKCoqBl8tJunAYtphDEKxQDZIV-M982TmmePGFwYAEJkTVAxxhFeZZeCGKxmFsMB3hguD73jJ8XzfD9RWwbZvyOBoU2AMEn2cHDMKQAB2T8Pzw0JCQcOCoBRRAJRGFDiMojDX2wY1OJo+8YCfbiDm-HZCDZXDf2EbQEh0PQUkMe91CfboxLfXZRPE1k+IIhkzm2ZhnHZHN5m5JBC35BBZTlcVJWBP5kyaYBUMZFTxJNDSdmYQDFgI5zmNU7ovJ2bpX0C7zIl85SQs8vYQu2MLtMcvyEDimK1NirzEpiZKNnUsK31ZGKsvkQkXkZXKMr2CqdnyQtjHMnIN1LLdWIIx8XOi7BmFCzruJoqB-3SIRIOglLhgcmJoA6rjuuC0LMskqBgI2ciTQm+QQGUqx4qC1KKv6igEDwfAkBQRStpwvLtoKuqfMcgTyvc3ZhMK57+pSRAOi6Xp1qcp8Nle6KAbZBLJLYx6Dmet93Owvpc1+9rNhhzNoYOLqf3hgbMeGqDmPG+9dJS66DL2bN+vzHk+RyXF8R+8HNle1TgYPbrir+8rCvRqAvO2dG2eSnmWdfKxuKF-mnxFtkufIoK+bByLGRlvYufyTC5fh+9Sol0Wuclg8+s5BrECalcEafQWPwOC3sP6wbyhx0bkIIqaWGw0UrCtt31fSYDnDW+9NsZVWZtfYO5oOnBqGYAgqF+pSg8w7porDuK2cJ-J3ONbMOUximzKpxBe2JX6ABlzeIWLsCtjYQar-r2FmJjuFQG07piIjGWYCudirlkeLexaMDo3IGKbp2kol-ve+2+K6-l+6JZRqGlZ74z7ekhBNFk-QFIVxAV9cqAV+9iLHPj-fduig-trTiWCm5wyycH0yjYQF04QRH5PXs+9Xg8AArKAORdioFwNyUwsJEDFkeGADAaBgwgG7NoJEQgMD2SgEIZQOA55CHKqtNarhgjlnrISQk+B6zsGvBWdA2huYghbtodQgRf4AKAcxEYQDwGgKwQgSErx2CvHAPA5QiCfgADF0GYLECdaouDpESEIYgLEjYrT8MEWAYRojawNn1B6KUJwuzoJtEw5Av0-6AOAaArhQCeFiKHAAR3rLCEA3gfgWkkTwqgYFTByOQKBbxoRgiQhtHaWADinEiO8AwkxZ0CK0MgSlPMDcEBwxwBwPwKYYiUNGnmQ6485BZLYcgBiOTQh5N+hEbQcSbHMRcBgZJqTTFQGSSgTJ8hsnMTKQ4E0+T5CVJkMACI4AIAyCAA}
\def\linkI{https://www.falstad.com/circuit/circuitjs.html?ctz=DwYwlgTgBAZgvAIgIwKgFwM6KQUwLQAMAbKmCNgQHQBMqIARogKwGoAOjCLqAbhM6gC2mbPgIB2BAD4AUFCjAAJoKgA7QYgBeOVTgh5qAFgBqqeAlZQwWZJXFJDAZibUAHEgCcbgq9+OxklAQNpaqiLRQ9DzhhggA9LLywIIq6ogANgD2AOZkeOqZGHhEBKawiJY8aNiUTKj0ONioGJyWihgAnhX1ik2R1RbN9L2DCXIKAB5QE4iGrlCO1FAEy2bdUBhg4ZZojQgAkgBcUAAiOJvZqugAhhgA1lAAFPQAlMdGUMZQcVAsUIIAQQ2mXSOkyYEUUB413SPBwUGuaCgSBWAAlNPFEpNprN5kRDMtVuVBhstshYug9gAVAAW8MclEcnygAHEAAr7ADyUEUEDAcIwCKgWVyIDwoLh6SgADk8ABZTkAZQAYgBRKnoGkQTIAV2yNKgACUAPosgDcRuNbJOUDuOBwbEFaDpUJhcKgmRgMCgAHcaWBQZr4WwwFdrAiw6o2DqkY8IOccGgXpRMeNgFMZgg5lBfITLOZLJtsBTdohlekOvRriAHoowJlFPCQFkMDhIfQOjycDBrjr0mgLZk2DooAByDtsW4YUfoTIIxSQ53wj5fbS6aCPHsYXbQDA+sBoEA0vCemDJz5u+GhxQ6kBoa6qEBXwXXDbpCF6FNjJKGnFZggrEgbhQIYAFQB4JDEpYPoFkI1wTNECBGKmSQ+n+oGAcBhj4sibhrKMWLAACKSJno6FgUB8wYci4iQeYjhCHshbkARaZSpm1FILRIFgUY+bhKg6QVIyKEKNAHG8aBPErI4BAUrBQS4qw34KGhEmAdxjiOEQNF0d0KnAGpiBaTpskEtRZn4cphEgH+Jm6Qs2kLHJ+ERCxuCEB4onpn+nhLHxeZWc0ZIOKgpYIBgIJghCrqwjg3kZtgXhQNQOErPxJJFuSYVMYyUCogANCBrgrJyNKCN5Rn-jJckgThllQd5kIcfVtXYQSfFBf8mTYNKhgOK43maOhOFcTp7UpaBXWCD1CBrnoBgmJVI0dVJLjzJ1jUGVV60ORNY1BQZgqZrtAW7QdjVQGw6wzfphGJVwhgdUgOnpV1WWhZSiATlO3nHbMo3cfteGXddJK3axSQPXMgFEPMb2XR9JZ7DAFZVjWy3qc55lgYsGXWWmVUWdQKwWdph2EekVxY3jdUEnj+HI4gJzed6mb2c4SwWS402zSKeQFEUJSmAZ2TkTJcPSdJFNpr+mYvTJLjIkBCw85dMECbA-DIJYgjwYhSCUEQdSRGAtw1Mb3nidgJSq-5Umc112sogTqF-gFnPY3bMtu+ztWexzauwQZ1tIQ7SsBU4EQKc7AFx95YuZpHiyTaT+I+wotl+-TStmSBixdSxTD4CiCd2W1Kd5+1Gc+fLT1SwjClIzliAmiyRUmtaCXuyVIEQYFiMhUzCAAMKmoq3nXIoABWUB7JBOBkpYOiIPsqg3neD5Pk8qIvEIYQWJQBAoP8SfNAMxiXlY6+3vej7xcKORNAZU+z3s1CBIv6wrwgyrxgAjjqHQIBOyPHRHvf4B8T6CHlgBc+iBFT7kPP6VQYsYAAKAY+Dogkn7IG8nOH+yENhg2oEwPSSF2DdCxEkNgUJ1gYCLq7BQcRMiEQIQfE2GASFkPwifMGKBqEKFoYhBiGwWJUAIHUAyLDZDADiOACAsggA}
\def\linkJ{https://www.falstad.com/circuit/circuitjs.html?ctz=DwYwlgTgBAZgvAIgIwKgFwM6KQUwLRIpRgjYAMAdAEyogBGiArGagA4MLOoBuETqAW0zZ8hBAD4AUFCjABAqADsBiAOaswAezzLNGVPAQso3NIkoA2VHRzmKZJDSgYOxgCYYAnuetvs1syNUFz8ggHopGTkFZTUNbVYBPQMfE0DLa1sjezIADmDXVA9vIKg6UKI6dILQlgjpWQAPKEbEABZcqABmKigyPpTSjDBEKmM0LIApAC4oABEcYdVFdABDDABrKAAKEABKWaR+gFkACQAvKABlAAUASSgQABtNEC3NbhxoVagjx4U0BBViAcOhNFAfjAnusABYIepRZqtBAdKAWNp9AawVLDbBtVATRAAFSBIL6FAsAHYlBgdn8QApNCsAGIAJTaewANFAIDgQWBPtBGFBWMyRU8AK60gCiVzmUDcWjciwoUAA8hK0KxNcQVmgYaDHFABABBZw4NBoMCKVQ7VkAfSZELomh1XSQeDaFnVMIEXPhkSaLXanVynX6xkMxlxyHx6CyGCeOAA7jzVhNycKAGphRQYbkYL5gRY8xZgDBoTTQLpdH0KdPoA1QY4AYQAqqqrk8wMqIBhZg6rgWk6mgRNOyBNKwcLN7gA5CGKNxQVkADR2VxbAGkIWhG6CobC9hQAw1gKzgwgAJzeqiMb2ou9WbGlZOGJwwPjIMg-3+UYwCKsjTcNgFKMNYYDrKBFjgQisipsiN5QE+vxUJ0T6DHUgbAA8yKOOh96-JSt73phwRPHY4FQDCIwIF0FBdKeUQwJeVBtP0YydF0ZAYpxgxEEkcRaHgiTJHBwAscibH9NxGKyd0PFkcamhCdouj6OJCGjOxyF5MhOkYS+WFnlpCDSchhHmbJZGaZe8lUGxCm8aRRlMfBdk8RZ3r2WGNnYbayLyT0MmeV0jAfj44kXoFnlIFevTyUgFiRqkb6pIBwGIPRjG2XhxFefpHEuVGbnAKZ9mETWJHPiV4lPCsgVdNV3RNchbGDHGhIIPOpXQI1zVtD0XlKV+NalWgTyXoNvQoRYYbDa5CpZBYFCUqIxhgConDGMmtE9KVAXtENaGdNN3RjH5Z6HQgc3ofNt3nSl4TYaZD0oZS4ULbVL2Xh9M2EbkSVfZFP3IoDzVhkVNUg2e9W-Z970I+1L6dVka6lSAl6QwV2NVUppCcKINDidduOtbjc2XVEy5g0DKHg788VKYJCDKjAqwSk8aCldFiAM3FvT88lSlpaUGUgXRDGlTTfNA8FUAM4ZJXKYgbMc1zB1Y3LYwK3L4VU0GyKUq17ERliysxkgqOIIeGAwtc266tq3PiasbgAFZQFkSCUqgOC0cYOCKIgrJlhWqyKGS2yaL6GB7IIwfZAJeE-sEgRXEWJa8sMFZViu9pXKVbue1kVH+6kQeIF2KYrumoLbDmebx8aieUMn-DOIEGpajqiY12O9eN3HqAvKo-jieClfIM+GCsNg0NmWwPiBlErAmDiBOwSvshr1kyMS9Gm+lWEmjYZPidUL7zhz9P-FL8gEgNKv690cEBP-lvT8717oxXjwWVvzsGQT+UQT5SGAGEcAEApBAA}


% numbered tasks: P = preparation (paper), S = simulation, D = discussion
\newcounter{task}[section]
\renewcommand{\thetask}{\thesection.\arabic{task}}
\newcommand{\task}[1]{\par\medskip\refstepcounter{task}%
  {\sffamily\bfseries\color{rwuviolet}Task \thetask}\ {\sffamily\small\color{rwucyandark}[#1]}\quad}
\newtcolorbox{modelbox}[1]{colback=rwucyan40!25,colframe=rwucyan,fonttitle=\sffamily\bfseries,title=#1,boxrule=0.6pt,arc=1mm}
\newtcolorbox{solbox}{colback=rwuviolet!6,colframe=rwuviolet,fonttitle=\sffamily\bfseries,title=Solution,boxrule=0.6pt,arc=1mm,breakable}
\ifdefined\solutions
  \newenvironment{solution}{\begin{solbox}}{\end{solbox}}
\else
  \usepackage{comment}
  \excludecomment{solution}
\fi
% worked solutions: numbered steps, then the answer to the question
\newlist{steps}{enumerate}{1}
\setlist[steps]{label={\sffamily\bfseries\color{rwuviolet}Step \arabic*:},leftmargin=*,align=left,itemsep=0.6ex}
\newcommand{\answer}{\par\smallskip{\sffamily\bfseries\color{rwuviolet}Answer:}\ }
% value measured in the simulation: empty cell for students, filled in the solutions
\ifdefined\solutions
  \newcommand{\simv}[1]{\textcolor{rwuviolet}{#1}}
\else
  \newcommand{\simv}[1]{}
\fi
% table of given (calculated) values; simulation columns are filled in by the students
\newenvironment{results}[2]{\par\smallskip\begin{center}\small\sffamily
  {\color{rwucyandark}\bfseries #1}\\[1mm]\renewcommand{\arraystretch}{1.25}\begin{tabular}{#2}\toprule}
  {\bottomrule\end{tabular}\end{center}}
% "open circuit X" link with the file name
\newcommand{\circuitlink}[2]{\href{#1}{\color{rwucyandark}\sffamily\bfseries $\blacktriangleright$\ open in Falstad}\ \ (\texttt{#2})}
\newcommand{\schematic}[2][0.62]{\begin{center}\includegraphics[width=#1\textwidth]{figures/falstad/#2_schematic}\end{center}}

\begin{document}

\begin{center}
  \rwulogo\\[6mm]
  {\sffamily\bfseries\LARGE\color{rwuviolet} GPIO Output Stages in Simulation}\\[2mm]
  {\sffamily\Large\color{rwucyandark} Push-pull, open drain and open collector: switching speed, voltage and current limits}\\[3mm]
  {\sffamily Embedded Computing -- Master of Engineering -- Worksheet\ifdefined\solutions\ with solutions\fi}
\end{center}

\section*{Overview}

Every GPIO pin ends in an output stage: two transistors (push-pull) or one transistor to ground (open drain,
open collector). In this worksheet you will take these stages apart in a circuit simulator. You will find out
how much current a pin can deliver, how fast it can switch a load, why a pull-up resistor makes a bus slow,
and which voltages and currents destroy a pin or the circuit behind it.

\textbf{Learning outcomes.} After this worksheet you can
\begin{itemize}
  \item model a GPIO output stage by on-resistance and saturation current and relate the model to the
        $V_{OH}$/$V_{OL}$ specifications of a data sheet,
  \item predict rise and fall times, the maximum toggle frequency and the dynamic power of a push-pull output,
  \item dimension the pull-up resistor of an open-drain bus such as I$^2$C from rise-time and $V_{OL}$ limits,
  \item judge where the maximum output voltage and current of a pin come from (body diodes, injection current,
        back-powering, transistor breakdown and saturation) and protect a driver against inductive kick-back,
  \item recognise when a PCB trace has to be treated as a transmission line and terminate it.
\end{itemize}

\textbf{Organisation.} About 4 hours, in teams of two. You do not have to calculate: all results of the
model are given in tables. Tasks marked \textsf{[P]} (preparation) ask you to read these values and understand
where they come from \emph{before} you open the simulation. \textsf{[S]} means simulate, measure and enter your
values in the empty column \emph{Simulation} of the tables. \textsf{[D]} means discuss and explain. Hand in a short report (PDF): for each task the completed tables, screenshots of the
scopes and your explanation. When a measurement differs from the model value by
more than about \SI{10}{\percent}, explain why.

\textbf{Tool.} Falstad CircuitJS runs in the browser: \url{https://www.falstad.com/circuit/}. Each part of this
sheet has a link that opens the prepared circuit. The circuit files are also in the folder
\href{https://github.com/mopfeil/EmbeddedComputing/tree/main/examples/interfaces/falstad}{\texttt{examples/interfaces/falstad/}}
of the lecture repository (open them with \emph{File $\rightarrow$ Open File}).

\begin{modelbox}{Working with CircuitJS}
\begin{itemize}[leftmargin=*]
  \item Hover over an element to see its voltages, currents and power in the lower right corner.
  \item Click on a logic input (\textsf{H}/\textsf{L}) or a switch to toggle it. Double-click an element to edit
        its value. Sliders on the right change the values marked in the circuit.
  \item Scopes: right-click $\rightarrow$ \emph{Properties} sets the time base (``scope speed''), the scale and
        shown values (e.g.\ average, peak-to-peak). Drag with the mouse in a scope to measure time differences.
  \item \emph{Reset} restarts the simulation, \emph{Run/STOP} freezes it for reading values.
  \item \emph{Options $\rightarrow$ Other Options} lets you change the simulation time step. If edges look
        jagged, make the time step smaller.
\end{itemize}
\end{modelbox}

\begin{modelbox}{The transistor model used in all circuits}
The MOSFETs follow the square-law (Shichman--Hodges) model with threshold $V_T=\SI{0.6}{V}$ and
\[
  I_D=\begin{cases}
     \beta\left[(V_{GS}-V_T)\,V_{DS}-\tfrac12 V_{DS}^2\right] & V_{DS} < V_{GS}-V_T \quad\text{(triode)}\\[1mm]
     \tfrac{\beta}{2}\,(V_{GS}-V_T)^2 & V_{DS} \ge V_{GS}-V_T \quad\text{(saturation)}
  \end{cases}
\]
with $\beta_N=\SI{12}{mA/V^2}$ for the NMOS and $\beta_P=\SI{8}{mA/V^2}$ for the PMOS of the output stage (PMOS:
use $V_{SG}$, $V_{SD}$). Both have a body diode. The supply is $V_{DD}=\SI{3.3}{V}$.
The model knows no gate or junction capacitances, no temperature and no breakdown. Pin and load capacitances are
drawn as explicit capacitors. Keep in mind what the model cannot show when you interpret the results.
\end{modelbox}

%======================================================================
\needspace{0.45\textheight}
\section{Push-pull output: static characteristic}
%======================================================================
\circuitlink{\linkA}{A\_pushpull\_characteristic.txt}

The logic input is the bit in the output register. Two inverters (the predriver and the CMOS stage) make the pin
follow it. A triangle voltage source acts as a curve tracer: it forces the pin voltage from \SI{-0.5}{V} to
\SI{3.8}{V}, and an ammeter measures the pin current (positive = current out of the pin). The left scope shows the
current over the voltage, the right scope the power in both transistors.

\schematic[0.55]{A}

\task{P} Table~1.1 gives the on-resistance (valid for small $V_{DS}$) and the saturation current of both
transistors for $V_{GS}=V_{DD}$. Which one is the stronger transistor, and why do chip designers make the PMOS
larger than the NMOS?

\begin{results}{Table 1.1: model values of the output transistors ($V_{GS}=\SI{3.3}{V}$)}{lccc}
  & formula & NMOS & PMOS\\ \midrule
  $\beta$ & & \SI{12}{mA/V^2} & \SI{8}{mA/V^2}\\
  $R_{on}$ & $1/(\beta(V_{GS}-V_T))$ & \SI{30.9}{\ohm} & \SI{46.3}{\ohm}\\
  $I_{sat}$ & $\tfrac{\beta}{2}(V_{GS}-V_T)^2$ & \SI{43.7}{mA} & \SI{29.2}{mA}\\
\end{results}

\begin{solution}
\begin{steps}
  \item Gate overdrive: with $V_{GS}=V_{DD}=\SI{3.3}{V}$ and $V_T=\SI{0.6}{V}$ both transistors have
        $V_{GS}-V_T=\SI{2.7}{V}$.
  \item On-resistance: for small $V_{DS}$ the triode equation reduces to $I_D\approx\beta(V_{GS}-V_T)V_{DS}$, so
        $R_{on}=V_{DS}/I_D=1/(\beta(V_{GS}-V_T))$:\\
        $R_{on,N}=1/(\SI{12}{mA/V^2}\cdot\SI{2.7}{V})=\SI{30.9}{\ohm}$,\quad
        $R_{on,P}=1/(\SI{8}{mA/V^2}\cdot\SI{2.7}{V})=\SI{46.3}{\ohm}$.
  \item Saturation current: $I_{sat}=\tfrac{\beta}{2}(V_{GS}-V_T)^2$ with $(\SI{2.7}{V})^2=\SI{7.29}{V^2}$:\\
        $I_{sat,N}=\SI{6}{mA/V^2}\cdot\SI{7.29}{V^2}=\SI{43.7}{mA}$,\quad
        $I_{sat,P}=\SI{4}{mA/V^2}\cdot\SI{7.29}{V^2}=\SI{29.2}{mA}$.
  \item Comparison: the NMOS has the lower $R_{on}$ and the higher $I_{sat}$, both by the factor
        $\beta_N/\beta_P=1.5$.
  \item Reason: $\beta=\mu C_{ox}W/L$. The mobility of holes (PMOS) is about 2--3 times lower than that of
        electrons (NMOS), so a PMOS needs a 2--3 times wider channel $W$ for the same $\beta$. Here the PMOS is
        larger, but not enough to fully compensate, which saves chip area.
\end{steps}
\answer The NMOS is the stronger transistor. Designers make the PMOS wider to compensate for the lower hole
mobility. It often stays somewhat weaker, which is why data sheets allow a larger $V_{OH}$ drop than $V_{OL}$ rise,
or specify asymmetric currents.
\end{solution}

\task{S} Set the output to \textsf{H} and then to \textsf{L}. From the curve, read $V_{OH}$ and $V_{OL}$ at the
currents of Table~1.2 (Run/STOP and hovering over the PIN label helps) and enter them next to the model values.
Also read the short-circuit currents (pin at \SI{0}{V} for \textsf{H}, pin at \SI{3.3}{V} for \textsf{L}).

\begin{results}{Table 1.2: output levels over the load current}{lcccccc}
  $|I|$ & \SI{2}{mA} & \SI{4}{mA} & \SI{8}{mA} & \SI{12}{mA} & \SI{20}{mA} & short circuit\\ \midrule
  $V_{OH}$, model & \SI{3.21}{V} & \SI{3.11}{V} & \SI{2.90}{V} & \SI{2.67}{V} & \SI{2.11}{V} & \SI{29.2}{mA}\\
  $V_{OH}$, simulation & \simv{3.21} & \simv{3.11} & \simv{2.90} & \simv{2.67} & \simv{2.11} & \simv{29.2}\\ \midrule
  $V_{OL}$, model & \SI{0.06}{V} & \SI{0.13}{V} & \SI{0.26}{V} & \SI{0.40}{V} & \SI{0.71}{V} & \SI{43.7}{mA}\\
  $V_{OL}$, simulation & \simv{0.06} & \simv{0.13} & \simv{0.26} & \simv{0.40} & \simv{0.71} & \simv{43.7}\\
\end{results}

\begin{solution}
\begin{steps}
  \item Output \textsf{H}: the PMOS conducts, the pin current flows out of the pin and the voltage across the
        PMOS is $V_{SD}=V_{DD}-V_{OH}$. For small currents the PMOS is in the triode region:
        $I=\beta_P\left[2.7\,V_{SD}-\tfrac12V_{SD}^2\right]$.
  \item Solve the quadratic equation for $V_{SD}$: $V_{SD}=2.7\,\si{V}-\sqrt{(2.7\,\si{V})^2-2I/\beta}$.
        Example $I=\SI{8}{mA}$: $2I/\beta_P=\SI{2}{V^2}$, $V_{SD}=2.7-\sqrt{5.29}=\SI{0.40}{V}$, so
        $V_{OH}=\SI{2.90}{V}$.
  \item Output \textsf{L}: the same formula with $\beta_N$ gives $V_{OL}=V_{DS}$. Example $I=\SI{8}{mA}$:
        $2I/\beta_N=\SI{1.33}{V^2}$, $V_{OL}=2.7-\sqrt{5.96}=\SI{0.26}{V}$.
  \item In the simulation: set the logic input, press Run/STOP while the sweep passes the desired current and
        hover over PIN. The measured values are:\\[1mm]
\begin{tabular}{lccccc}
  \toprule
  $|I|$ / mA & 2 & 4 & 8 & 12 & 20\\
  \midrule
  $V_{OH}$ / V & 3.21 & 3.11 & 2.90 & 2.67 & 2.11\\
  $V_{OL}$ / V & 0.06 & 0.13 & 0.26 & 0.40 & 0.71\\
  \bottomrule
\end{tabular}
  \item Short circuit: for \textsf{H} with the pin at \SI{0}{V} the PMOS sees $V_{SD}=\SI{3.3}{V}>V_{GS}-V_T$, so it
        is saturated and delivers $I_{sat,P}=\SI{29.2}{mA}$. For \textsf{L} with the pin at \SI{3.3}{V} the NMOS
        delivers $I_{sat,N}=\SI{43.7}{mA}$.
\end{steps}
\answer Model and simulation agree to within the reading accuracy. Short-circuit currents: \SI{29.2}{mA}
(high, PMOS saturated) and \SI{43.7}{mA} (low, NMOS saturated).
\end{solution}

\task{D} The RP2040 data sheet specifies $V_{OH}\ge\SI{2.62}{V}$ and $V_{OL}\le\SI{0.5}{V}$ at a load current
equal to the selected drive strength (2, 4, 8 or \SI{12}{mA}) and $IOVDD=\SI{3.3}{V}$. The model delivers
\SI{12.8}{mA} at $V_{OH}=\SI{2.62}{V}$ and sinks \SI{12}{mA} at $V_{OL}=\SI{0.40}{V}$. Which drive strength setting
does the model represent? Drive strength is implemented by switching transistor ``fingers'' in parallel, i.e.\ by
scaling $\beta$ ($I_D\propto\beta$): the \SI{4}{mA} setting needs $\beta_P=\SI{2.5}{mA/V^2}$. What fraction of
the fingers is that, and what does it mean for the switching speed (part 3)?

\begin{solution}
\begin{steps}
  \item High side: at $V_{OH}=\SI{2.62}{V}$ the PMOS has $V_{SD}=\SI{0.68}{V}$ and delivers
        $I=\SI{8}{mA/V^2}\cdot(2.7\cdot0.68-\tfrac12\cdot0.68^2)\,\si{V^2}=\SI{12.8}{mA}\ge\SI{12}{mA}$.
  \item Low side: at $V_{OL}=\SI{0.4}{V}$ the NMOS sinks
        $I=\SI{12}{mA/V^2}\cdot(2.7\cdot0.4-\tfrac12\cdot0.4^2)\,\si{V^2}=\SI{12}{mA}$, exactly the required value.
  \item Both data-sheet conditions are met at \SI{12}{mA} (and not at a higher current), so the model represents
        the \SI{12}{mA} setting.
  \item \SI{4}{mA} setting: at the fixed operating point $V_{SD}=\SI{0.68}{V}$ the current scales with $\beta$:
        $\beta_P=\SI{8}{mA/V^2}\cdot\SI{4}{mA}/\SI{12.8}{mA}=\SI{2.5}{mA/V^2}$. That is $2.5/8\approx 1/3$ of
        the fingers.
  \item Switching speed: with one third of $\beta$, $R_{on}$ is three times larger and $I_{sat}$ three times
        smaller. Edge times are $\propto C_L/\beta$ (part 3), so the pin switches about three times more slowly.
\end{steps}
\answer The model represents the \SI{12}{mA} setting. The \SI{4}{mA} setting uses about one third of the
fingers and makes rise and fall times about three times longer for the same load capacitance.
\end{solution}

\task{S/D} A red LED ($V_F\approx\SI{2.0}{V}$) with a \SI{100}{\ohm} series resistor is connected from the pin to
GND. The naive calculation gives $I=(\SI{3.3}{V}-\SI{2.0}{V})/\SI{100}{\ohm}=\SI{13}{mA}$. Draw the load line
$I=(V_{PIN}-V_F)/R$ into your measured characteristic (it passes through \SI{2.0}{V}/\SI{0}{mA} and
\SI{3.3}{V}/\SI{13}{mA}) and read the operating point. The model gives $V_{PIN}=\SI{2.86}{V}$ and $I=\SI{8.6}{mA}$.
Why is the current a third lower than the naive value? Optional: verify it by replacing the sweep source with an
LED and a resistor (\emph{Draw $\rightarrow$ Outputs and Labels $\rightarrow$ Add LED}).

\begin{solution}
\begin{steps}
  \item Load line: the LED branch draws $I=(V_{PIN}-\SI{2.0}{V})/\SI{100}{\ohm}$. Draw the straight line through
        \SI{2.0}{V}/\SI{0}{mA} and \SI{3.3}{V}/\SI{13}{mA} into the \textsf{H} curve.
  \item Operating point: the pin current must satisfy both the output characteristic and the load line. With
        $V_{SD}=V_{DD}-V_{PIN}$: $\SI{10}{mA/V}\cdot(1.3\,\si{V}-V_{SD})=\SI{8}{mA/V^2}\,(2.7\,V_{SD}-\tfrac12V_{SD}^2)$
        $\Rightarrow 4V_{SD}^2-31.6\,V_{SD}+13=0\Rightarrow V_{SD}=\SI{0.44}{V}$.
  \item Result: $V_{PIN}=\SI{2.86}{V}$, $I=(2.86-2.0)\,\si{V}/\SI{100}{\ohm}=\SI{8.6}{mA}$. The graphical
        intersection in the simulation gives the same values.
  \item Interpretation: at this point the PMOS behaves like a resistor of
        $\SI{0.44}{V}/\SI{8.6}{mA}\approx\SI{50}{\ohm}$, which is in series with the \SI{100}{\ohm}. The total
        resistance is \SI{150}{\ohm} instead of \SI{100}{\ohm}, so the current drops to about two thirds.
\end{steps}
\answer The output resistance of the pin (about \SI{50}{\ohm}) is not negligible compared to \SI{100}{\ohm}.
For a defined LED current either choose a larger series resistor (then the pin resistance becomes negligible) or
include the output stage in the calculation.
\end{solution}

\task{D} Look at the curve below \SI{0}{V} and above \SI{3.3}{V}. Which element conducts there? Relate this to the
absolute maximum rating ``pin voltage $-0.5\ldots IOVDD+\SI{0.5}{V}$'' of the RP2040 and to the term
\emph{injection current}. Why can a large injection current upset an analogue measurement or even a digital
circuit on the same chip?

\begin{solution}
\begin{steps}
  \item Above $V_{DD}$: the drain of the PMOS (the pin) becomes more positive than its source and body
        ($V_{DD}$). The body diode drain$\rightarrow$body is forward biased and conducts current from the pin
        into the supply. In a real chip the ESD clamp diode to $V_{DD}$ is in parallel.
  \item Below \SI{0}{V}: the body diode of the NMOS (body at GND, drain at the pin) is forward biased and
        conducts current from GND out of the pin.
  \item Shape of the curve: from about \SI{0.4}{V} beyond the rails the diode current rises exponentially
        (in the model to about \SI{9}{mA} at \SI{3.7}{V}).
  \item Absolute maximum rating: the limits $-0.5\ldots IOVDD+\SI{0.5}{V}$ keep these diodes almost off. The
        current that still flows through them is called injection current.
  \item Consequences: injected current flows as minority carriers into the substrate or into the supply rail.
        It can shift the ADC reference and neighbouring analogue channels, trigger the parasitic thyristor of
        the CMOS structure (latch-up, a short circuit between $V_{DD}$ and GND), or raise $V_{DD}$ (see
        part~\ref{sec:overvoltage}).
\end{steps}
\answer The body diodes (and ESD diodes) conduct outside the supply range. The absolute maximum rating keeps
them off; a large injection current disturbs analogue measurements and can latch up the chip.
\end{solution}

\task{S/D} Read the maximum power in the transistors from the right scope (model: \SI{96}{mW} in the PMOS for a
short circuit to GND, \SI{110}{mW} at \SI{-0.5}{V}). Many microcontrollers allow \SIrange{50}{100}{mA} in total for
all pins of one supply domain. What does this mean for the number of pins that may drive LEDs at the same time,
and what happens to the chip temperature if a pin is shorted to GND?

\begin{solution}
\begin{steps}
  \item Short circuit to GND: the PMOS is saturated, carries \SI{29.2}{mA} and has the full supply across it:
        $P=\SI{3.3}{V}\cdot\SI{29.2}{mA}\approx\SI{96}{mW}$.
  \item Pin at \SI{-0.5}{V}: the current stays at $I_{sat}$, but the voltage rises to \SI{3.8}{V}:
        $P\approx\SI{3.8}{V}\cdot\SI{29.2}{mA}\approx\SI{110}{mW}$.
  \item Number of LEDs: each LED from task 1.4 draws \SI{8.6}{mA}. With a budget of \SI{50}{mA} per supply
        domain: $\SI{50}{mA}/\SI{8.6}{mA}=5.8$, so at most five LEDs at the same time (with \SI{100}{mA}: eleven).
  \item Temperature: the \SIrange{96}{110}{mW} are dissipated in a tiny transistor area at the pad. The chip
        as a whole heats up only a little, but locally the temperature rises strongly. The current also flows
        through thin supply lines on the chip (electromigration).
\end{steps}
\answer Only a few pins may drive LEDs at the same time (about five at \SI{50}{mA}). A single short-circuited
pin usually survives, but it heats up locally and ages. This is why per-pin and per-domain limits exist, and
larger loads need an external driver (part~\ref{sec:oc}).
\end{solution}

%======================================================================
\needspace{0.45\textheight}
\section{Two push-pull outputs on one line}
%======================================================================
\circuitlink{\linkC}{C\_bus\_contention.txt}

\schematic[0.62]{C}

\task{P/S} Output A drives high, output B drives low. The line voltage follows from $I_{D,P,A}(V)=I_{D,N,B}(V)$;
both transistors work in the triode region. Measure the values of Table~2.1 in the simulation.

\begin{results}{Table 2.1: contention between two push-pull outputs}{lcccc}
  & line voltage & current & power PMOS A & power NMOS B\\ \midrule
  model & \SI{1.09}{V} & \SI{28.2}{mA} & \SI{62}{mW} & \SI{31}{mW}\\
  simulation & \simv{1.09} & \simv{28.2} & \simv{62.3} & \simv{30.8}\\
\end{results}

\begin{solution}
\begin{steps}
  \item Current path: from $V_{DD}$ through the PMOS of A, along the line, through the NMOS of B to GND. Both
        transistors carry the same current $I$.
  \item Equation: the PMOS of A sees $V_{SD}=3.3\,\si{V}-V$, the NMOS of B sees $V_{DS}=V$. Assuming the triode
        region for both:
        $\beta_P\left[2.7\,(3.3-V)-\tfrac12(3.3-V)^2\right]=\beta_N\left[2.7\,V-\tfrac12V^2\right]$.
  \item Solving (numerically, or by trying values) gives $V=\SI{1.09}{V}$. Check: NMOS
        $\SI{12}{mA/V^2}\cdot(2.7\cdot1.09-\tfrac12\cdot1.09^2)\,\si{V^2}=\SI{28.2}{mA}$; PMOS with
        $V_{SD}=\SI{2.21}{V}$: $\SI{8}{mA/V^2}\cdot(2.7\cdot2.21-\tfrac12\cdot2.21^2)\,\si{V^2}=\SI{28.2}{mA}$.
  \item Check the assumption: $V_{DS}=\SI{1.09}{V}$ and $V_{SD}=\SI{2.21}{V}$ are both below
        $V_{GS}-V_T=\SI{2.7}{V}$, so both transistors are indeed in the triode region.
  \item Power: $P_{P,A}=\SI{2.21}{V}\cdot\SI{28.2}{mA}=\SI{62}{mW}$,
        $P_{N,B}=\SI{1.09}{V}\cdot\SI{28.2}{mA}=\SI{31}{mW}$. The weaker PMOS takes the larger voltage and
        twice the power.
  \item Simulation: \SI{1.09}{V}, \SI{28.20}{mA}, \SI{62.3}{mW} and \SI{30.8}{mW}.
\end{steps}
\answer The line settles at \SI{1.09}{V} with \SI{28.2}{mA} contention current; model and simulation agree.
\end{solution}

\task{D} What logic level does a receiver see on the line?

\begin{solution}
\begin{steps}
  \item Input thresholds of a \SI{3.3}{V} CMOS input: $V_{IL,max}\approx\SI{0.8}{V}$ (LVTTL) or
        $0.3\,V_{DD}=\SI{1.0}{V}$, $V_{IH,min}\approx\SI{2.0}{V}$ or $0.7\,V_{DD}=\SI{2.3}{V}$.
  \item The line voltage \SI{1.09}{V} is above $V_{IL,max}$ and below $V_{IH,min}$.
  \item In this range the receiver may read 0 or 1, the value can differ between receivers and change with
        temperature, and the input stage draws cross current (task 3.5).
\end{steps}
\answer The level is undefined: \SI{1.09}{V} lies in the forbidden zone between $V_{IL,max}$ and
$V_{IH,min}$.
\end{solution}

\task{D} Name three situations in embedded practice where this happens. How do tri-state outputs, open-drain
outputs and series resistors prevent the damage? (A series resistor of about \SI{290}{\ohm} in each output
limits the contention current to \SI{5}{mA}: $2R+R_{on,P}+R_{on,N}=\SI{3.3}{V}/\SI{5}{mA}$. What does it cost?)

\begin{solution}
\begin{steps}
  \item Situations (any three): two masters on a shared line (bus arbitration error); a bidirectional data line
        where both sides drive at the same time during the turn-around; bit-banged I$^2$C with the pin
        configured push-pull; a debugger and the board both driving reset; two boards connected TX to TX; a pin
        shorted to its neighbour by a solder bridge; a pin configured as output by mistake while an external
        signal drives it.
  \item Tri-state outputs: each driver can switch to high impedance. Only one driver is enabled at a time;
        the protocol (or the firmware) must guarantee this, so a bug still causes contention.
  \item Open-drain outputs: no output drives high, only the pull-up does. If two outputs disagree, the low
        output wins and the current is limited by the pull-up resistor. Contention is harmless by design.
  \item Series resistors: the contention current flows through $2R+R_{on,P}+R_{on,N}$. For \SI{5}{mA}:
        $2R+\SI{46.3}{\ohm}+\SI{30.9}{\ohm}=\SI{3.3}{V}/\SI{5}{mA}=\SI{660}{\ohm}\Rightarrow R\approx\SI{290}{\ohm}$
        per output (standard value \SI{330}{\ohm}).
  \item Cost: \SI{330}{\ohm} in series with the load capacitance increases the time constant about tenfold
        (compared to $R_{on}\approx\SI{30}{\ohm}\ldots\SI{46}{\ohm}$), and the resistor drops voltage under DC load.
\end{steps}
\answer Contention happens whenever two drivers may be active on one line. Tri-state prevents it by protocol,
open drain by design, series resistors by limiting the current, at the cost of slower edges (part 3).
\end{solution}

%======================================================================
\needspace{0.45\textheight}
\section{Push-pull output: switching speed and dynamic power}
%======================================================================
\circuitlink{\linkB}{B\_pushpull\_capacitive\_load.txt}

The square wave is the program toggling the pin. The predriver is a logic inverter with an adjustable slew rate
(how fast the gate voltage changes). The output drives the load capacitance $C_L$ (pin, trace, inputs of other
chips). An ammeter measures the current taken from the supply.

\schematic[0.55]{B}

\task{P} Table~3.1 gives the rise and fall times (\SIrange{10}{90}{\percent}) from two models: (a) the RC
approximation $t_r=2.2\,R_{on}C_L$, (b) the square-law model, in which the transistor starts in saturation
(constant current $I_{sat}$) and enters the triode region when $V_{DS}=V_{GS}-V_T$ (integration of
$dt=C_L\,dV/I_D(V)$). Why is the RC approximation too optimistic?

\begin{solution}
\begin{steps}
  \item RC model (a): charging through a constant resistance gives $V(t)=V_{DD}(1-e^{-t/\tau})$; from
        \SI{10}{\percent} to \SI{90}{\percent} it takes $\tau\ln 9=2.2\,\tau$ with $\tau=R_{on}C_L$.
        Rising edge (PMOS): $t_r=2.2\cdot\SI{46.3}{\ohm}\cdot C_L$, falling edge (NMOS):
        $t_f=2.2\cdot\SI{30.9}{\ohm}\cdot C_L$. For \SI{10}{pF}, \SI{100}{pF}, \SI{1}{nF}:
        \SI{1.0}{ns}/\SI{0.68}{ns}, \SI{10.2}{ns}/\SI{6.8}{ns}, \SI{102}{ns}/\SI{68}{ns}.
  \item Square-law model (b), rising edge: the pin starts at \SI{0}{V}, so the PMOS starts with
        $V_{SD}=\SI{3.3}{V}>a=V_{GS}-V_T=\SI{2.7}{V}$, i.e.\ in saturation. It stays saturated until
        $V_{SD}=\SI{2.7}{V}$, i.e.\ until the pin reaches \SI{0.6}{V}.
  \item Saturation phase: constant current, so the pin voltage rises linearly. From the \SI{10}{\percent} point
        (\SI{0.33}{V}) to \SI{0.6}{V}: $t_1=C_L\cdot\SI{0.27}{V}/\SI{29.2}{mA}$ (\SI{0.92}{ns} for \SI{100}{pF}).
  \item Triode phase: $C_L\,dV=I_D\,dt$ with $x=V_{SD}$ falling from \SI{2.7}{V} to \SI{0.33}{V}
        (\SI{90}{\percent} point):
\[
  t_2=\frac{C_L}{\beta}\int_{0.33}^{2.7}\frac{dx}{x(a-x/2)}=\frac{C_L}{\beta a}\left[\ln\frac{x}{2a-x}\right]_{0.33}^{2.7}
     =\frac{C_L}{\beta}\cdot\SI{1.012}{V^{-1}}.
\]
        For \SI{100}{pF}: $t_2=\SI{100}{pF}/\SI{8}{mA/V^2}\cdot\SI{1.012}{V^{-1}}=\SI{12.65}{ns}$.
  \item Rise time: $t_r=t_1+t_2=\SI{13.6}{ns}$ for \SI{100}{pF}; it scales with $C_L$: \SI{1.36}{ns},
        \SI{13.6}{ns}, \SI{136}{ns}.
  \item Falling edge: the same calculation with the NMOS ($I_{sat}=\SI{43.7}{mA}$, $\beta=\SI{12}{mA/V^2}$):
        $t_1=\SI{0.62}{ns}$, $t_2=\SI{8.43}{ns}$ for \SI{100}{pF}, so $t_f=\SI{0.90}{ns}$, \SI{9.0}{ns},
        \SI{90}{ns}.
  \item Comparison: model (b) is about \SI{33}{\percent} slower than model (a), i.e.\ the RC estimate is about
        \SI{25}{\percent} too short.
\end{steps}
\answer The RC model assumes the transistor is a resistor from the very beginning, so at the start of the edge it
would deliver $V_{DD}/R_{on}=\SI{71}{mA}$ (PMOS). In reality the transistor is saturated and delivers only
$I_{sat}=\SI{29.2}{mA}$ until $V_{DS}$ has dropped to $V_{GS}-V_T$.
\end{solution}

\task{S} Set the predriver slew rate to its maximum and measure $t_r$ and $t_f$ for the three capacitances (set the
scope time base accordingly, or zoom in with Run/STOP). For which $C_L$ does the predriver dominate the edge?

\begin{results}{Table 3.1: rise and fall times at the pin}{lcccccc}
  & \multicolumn{2}{c}{$C_L=\SI{10}{pF}$} & \multicolumn{2}{c}{$C_L=\SI{100}{pF}$} & \multicolumn{2}{c}{$C_L=\SI{1}{nF}$}\\
  & $t_r$ & $t_f$ & $t_r$ & $t_f$ & $t_r$ & $t_f$\\ \midrule
  (a) RC model & \SI{1.0}{ns} & \SI{0.68}{ns} & \SI{10.2}{ns} & \SI{6.8}{ns} & \SI{102}{ns} & \SI{68}{ns}\\
  (b) square-law model & \SI{1.36}{ns} & \SI{0.90}{ns} & \SI{13.6}{ns} & \SI{9.0}{ns} & \SI{136}{ns} & \SI{90}{ns}\\
  simulation & \simv{1.5} & \simv{1.1} & \simv{13.6} & \simv{9.1} & \simv{136} & \simv{90}\\
\end{results}

\begin{solution}
\begin{steps}
  \item Set the slew rate slider to its maximum (\SI{2}{V/ns} or more) and the scope time base to a few edge
        times. Measure from \SI{0.33}{V} to \SI{2.97}{V} (\SIrange{10}{90}{\percent}).
  \item Measured values: \SI{10}{pF}: \SI{1.5}{ns}/\SI{1.1}{ns}; \SI{100}{pF}: \SI{13.6}{ns}/\SI{9.1}{ns};
        \SI{1}{nF}: \SI{136}{ns}/\SI{90}{ns}.
  \item Comparison: for \SI{100}{pF} and \SI{1}{nF} the simulation matches model (b) exactly. For \SI{10}{pF}
        it is about \SI{10}{\percent} (rising) and \SI{20}{\percent} (falling) slower than model (b).
  \item Reason: the gate voltage itself needs time to swing (\SI{3.3}{V} at \SI{2}{V/ns}: about \SI{1.5}{ns}).
        Model (b) assumes an instantaneous gate step. As long as the output edge is much longer than the gate
        edge, this does not matter.
\end{steps}
\answer At \SI{10}{pF} (and below) the predriver dominates the edge; from about \SI{100}{pF} on the output
transistors and $C_L$ determine it.
\end{solution}

\task{S/D} Increase the toggle frequency with $C_L=\SI{1}{nF}$ until the pin no longer reaches valid logic levels
($V_{IH}=0.7\,V_{DD}$, $V_{IL}=0.3\,V_{DD}$). Compare with the rule of thumb $f_{max}\approx 1/(2\,t_r)\approx
\SI{3.7}{MHz}$ (slower edge). Why do data sheets specify the maximum GPIO frequency only together with a load
capacitance?

\begin{solution}
\begin{steps}
  \item Thresholds: $V_{IH}=0.7\cdot\SI{3.3}{V}=\SI{2.31}{V}$, $V_{IL}=0.3\cdot\SI{3.3}{V}=\SI{0.99}{V}$. The
        pin must exceed \SI{2.31}{V} in every high phase and fall below \SI{0.99}{V} in every low phase.
  \item Rule of thumb: each half period must be at least as long as the slower edge, here $t_r=\SI{136}{ns}$:
        $f_{max}\approx1/(2\cdot\SI{136}{ns})=\SI{3.7}{MHz}$.
  \item Simulation: up to about \SI{3.7}{MHz} the pin still reaches almost the full swing. Above that the edges
        turn into exponential segments and the swing shrinks around $V_{DD}/2$. At about \SIrange{5}{6}{MHz}
        the levels no longer reach $V_{IH}$ and $V_{IL}$.
  \item The rule of thumb is slightly pessimistic, because valid levels need only \SI{70}{\percent} and
        \SI{30}{\percent} of the swing, not \SI{90}{\percent} and \SI{10}{\percent}.
\end{steps}
\answer The edge times are proportional to $C_L$, so the maximum frequency is inversely proportional to $C_L$.
A frequency without a load capacitance is meaningless; data sheets therefore state e.g.\ ``\SI{50}{MHz} at
\SI{30}{pF}''.
\end{solution}

\task{S} Set $f=\SI{1}{MHz}$ and the fastest slew rate. Read the average supply current (right scope) for the three
capacitances and compare with $I=C_LV_{DD}f$. The energy taken from the supply per cycle is $C_LV_{DD}^2$,
independent of $R_{on}$. Where is it dissipated?

\begin{results}{Table 3.2: average supply current at \SI{1}{MHz}}{lccc}
  & $C_L=\SI{10}{pF}$ & $C_L=\SI{100}{pF}$ & $C_L=\SI{1}{nF}$\\ \midrule
  $C_LV_{DD}f$ & \SI{33}{\micro A} & \SI{330}{\micro A} & \SI{3.3}{mA}\\
  simulation & \simv{34} & \simv{330} & \simv{3300}\\
\end{results}

\begin{solution}
\begin{steps}
  \item Charge per cycle: on every rising edge the supply delivers $Q=C_LV_{DD}$ to the load capacitance; on the
        falling edge this charge flows to GND, not back into the supply.
  \item Average current: $I=Q\cdot f=C_LV_{DD}f$, e.g.\ $\SI{1}{nF}\cdot\SI{3.3}{V}\cdot\SI{1}{MHz}=\SI{3.3}{mA}$;
        for \SI{10}{pF} and \SI{100}{pF}: \SI{33}{\micro A} and \SI{330}{\micro A}.
  \item Simulation: \SI{34}{\micro A}, \SI{330}{\micro A}, \SI{3300}{\micro A}. The small excess at \SI{10}{pF}
        is cross current during the gate edge (task 3.5).
  \item Energy: the supply delivers $E=QV_{DD}=C_LV_{DD}^2$ per cycle. At the end of the rising edge the
        capacitor stores $\tfrac12C_LV_{DD}^2$, so the other half has been converted to heat in the PMOS. This
        holds for any resistance: a smaller $R_{on}$ means a larger current for a shorter time.
  \item On the falling edge the stored $\tfrac12C_LV_{DD}^2$ is converted to heat in the NMOS.
\end{steps}
\answer The energy $C_LV_{DD}^2$ per cycle is dissipated half in the PMOS (charging) and half in the NMOS
(discharging). $P=C_LV_{DD}^2f$ is the dynamic power of all CMOS logic.
\end{solution}

\task{S/D} Set $C_L=\SI{1}{pF}$ and reduce the predriver slew rate to \SI{0.02}{V/ns}. Observe the supply current.
Both transistors conduct while $V_T<V_G<V_{DD}-V_T$ (\SIrange{0.6}{2.7}{V}), at \SI{0.02}{V/ns} that is
\SI{105}{ns} per edge. Measure the average supply current and compare it with $C_LV_{DD}f=\SI{3.3}{\micro A}$
(simulation: \underline{\hspace{2cm}}\simv{376\,µA}). Why do microcontroller input pins have Schmitt triggers, and
why does a slowly changing analogue voltage on a digital input increase the current consumption?

\begin{solution}
\begin{steps}
  \item Conduction window: the NMOS conducts for $V_G>V_T=\SI{0.6}{V}$, the PMOS for
        $V_G<V_{DD}-V_T=\SI{2.7}{V}$. Between \SI{0.6}{V} and \SI{2.7}{V} both are on.
  \item Duration: $\SI{2.1}{V}/\SI{0.02}{V/ns}=\SI{105}{ns}$ per edge, i.e.\ \SI{210}{ns} per period of
        \SI{1}{\micro s}.
  \item During this time a cross-conduction (shoot-through) current flows directly from $V_{DD}$ through both
        transistors to GND. In the simulation its peak is \SI{5.8}{mA} (at $V_G\approx V_{DD}/2$).
  \item Average supply current: with the fastest slew rate about \SI{6}{\micro A}, with \SI{0.02}{V/ns}
        \SI{376}{\micro A}, more than 100 times $C_LV_{DD}f=\SI{3.3}{\micro A}$. The cross current, not the load,
        determines the power.
  \item Transfer to inputs: the input buffer of a pin is also a CMOS inverter. A slow input edge or an analogue
        voltage near $V_{DD}/2$ keeps both of its transistors on, permanently for a constant voltage.
\end{steps}
\answer A Schmitt trigger switches quickly and only once even for a slow or noisy input edge (hysteresis), so the
cross current flows only briefly and the output does not oscillate. An analogue voltage on a digital input keeps
the input buffer in its conduction window and causes a static cross current; unused or analogue pins should
therefore have the digital input buffer disabled.
\end{solution}

%======================================================================
\needspace{0.45\textheight}
\section{Open-drain bus versus push-pull}
%======================================================================
\circuitlink{\linkD}{D\_open\_drain\_bus.txt}

Two open-drain devices share a line with one pull-up resistor $R_P$, as on an I$^2$C bus. Device~1 sends a clock;
device~2 can hold the line low. A push-pull stage drives the same capacitance for comparison.

\schematic[0.62]{D}

\task{P} The I$^2$C rise time is measured from \SI{30}{\percent} to \SI{70}{\percent} of $V_{DD}$. For the RC charge
$V(t)=V_{DD}(1-e^{-t/\tau})$ it is $t_r=\tau\ln(0.7/0.3)=0.8473\,R_PC_b$, i.e.\ \SI{796}{ns} for
$R_P=\SI{4.7}{k\ohm}$ and $C_b=\SI{200}{pF}$.

\begin{solution}
\begin{steps}
  \item When the open-drain transistor switches off, only $R_P$ charges the bus capacitance:
        $V(t)=V_{DD}(1-e^{-t/\tau})$ with $\tau=R_PC_b$.
  \item Time to \SI{30}{\percent}: $1-e^{-t_{30}/\tau}=0.3\Rightarrow t_{30}=\tau\ln(1/0.7)=0.357\,\tau$.
  \item Time to \SI{70}{\percent}: $1-e^{-t_{70}/\tau}=0.7\Rightarrow t_{70}=\tau\ln(1/0.3)=1.204\,\tau$.
  \item Rise time: $t_r=t_{70}-t_{30}=\tau\ln(0.7/0.3)=0.8473\,\tau$.
  \item Numbers: $\tau=\SI{4.7}{k\ohm}\cdot\SI{200}{pF}=\SI{940}{ns}$, $t_r=0.8473\cdot\SI{940}{ns}=\SI{796}{ns}$.
\end{steps}
\answer $t_r=0.8473\,R_PC_b=\SI{796}{ns}$, below the Standard-mode limit of \SI{1000}{ns}, but far above the
Fast-mode limit of \SI{300}{ns}.
\end{solution}

\task{S} Measure the rise and fall times of BUS and the rise time of PP (same thresholds). Explain the difference
between rising and falling edge of the open-drain line.

\begin{results}{Table 4.1: edges with $R_P=\SI{4.7}{k\ohm}$, $C_b=\SI{200}{pF}$ (\SIrange{30}{70}{\percent})}{lccc}
  & $t_r$ BUS & $t_f$ BUS & $t_r$ PP\\ \midrule
  model (RC) & $0.8473\,R_PC_b=\SI{796}{ns}$ & $0.8473\,R_{on,N}C_b\approx\SI{5}{ns}$ & $0.8473\,R_{on,P}C_b\approx\SI{8}{ns}$\\
  simulation & \simv{796} & \simv{8} & \simv{12}\\
\end{results}

\begin{solution}
\begin{steps}
  \item Measure with the thresholds \SI{0.99}{V} and \SI{2.31}{V} (\SIrange{30}{70}{\percent} of
        \SI{3.3}{V}): BUS $t_r=\SI{796}{ns}$, BUS $t_f=\SI{8}{ns}$, PP $t_r=\SI{12}{ns}$.
  \item The BUS rising edge matches the RC model exactly, because the pull-up is a real resistor.
  \item The falling edge of BUS and the rising edge of PP are somewhat slower than the RC estimate
        (\SI{5}{ns}, \SI{8}{ns}), because the transistors start in saturation (task 3.1).
  \item Falling edge of BUS: the NMOS actively discharges $C_b$ with about \SI{31}{\ohm}. Rising edge: no
        transistor drives high, only $R_P=\SI{4.7}{k\ohm}$ charges $C_b$.
  \item Ratio: $\SI{4.7}{k\ohm}/\SI{31}{\ohm}\approx150$; the edges differ by a factor of about 100 (796 versus
        \SI{8}{ns}).
\end{steps}
\answer The rising edge of an open-drain line is about 100 times slower than the falling edge, because it is
driven only by the pull-up resistor. This asymmetry is the price for the wired-AND function.
\end{solution}

\task{P} Table~4.2 gives the permitted range of $R_P$ for $C_b=\SI{200}{pF}$ and $V_{DD}=\SI{3.3}{V}$ according to
the I$^2$C specification (NXP UM10204): $V_{OL}\le\SI{0.4}{V}$ at a sink current of \SI{3}{mA}, $t_r\le\SI{1000}{ns}$
(Standard-mode, \SI{100}{kHz}) and $t_r\le\SI{300}{ns}$ (Fast-mode, \SI{400}{kHz}). Which limit comes from which
requirement?

\begin{results}{Table 4.2: pull-up range for $C_b=\SI{200}{pF}$}{lccc}
  & $R_{P,min}=\dfrac{V_{DD}-V_{OL}}{I_{OL}}$ & $R_{P,max}=\dfrac{t_r}{0.8473\,C_b}$ & standard values\\ \midrule
  Standard-mode & \SI{967}{\ohm} & \SI{5.9}{k\ohm} & \SI{1}{k\ohm}\ldots\SI{5.6}{k\ohm}\\
  Fast-mode & \SI{967}{\ohm} & \SI{1.77}{k\ohm} & \SI{1}{k\ohm}\ldots\SI{1.5}{k\ohm}\\
\end{results}

\begin{solution}
\begin{steps}
  \item Lower limit from the low level: while a device pulls the line low, the current through $R_P$ is
        $(V_{DD}-V_{OL})/R_P$. A device must only be able to sink \SI{3}{mA} at $V_{OL}\le\SI{0.4}{V}$, so
        $R_P\ge(\SI{3.3}{V}-\SI{0.4}{V})/\SI{3}{mA}=\SI{967}{\ohm}$. A smaller $R_P$ would demand more current
        than a compliant device guarantees.
  \item Upper limit from the rise time: $t_r=0.8473\,R_PC_b\le t_{r,max}\Rightarrow R_P\le t_{r,max}/(0.8473\,C_b)$.
  \item Standard-mode: $R_{P,max}=\SI{1000}{ns}/(0.8473\cdot\SI{200}{pF})=\SI{5.9}{k\ohm}$.
  \item Fast-mode: $R_{P,max}=\SI{300}{ns}/(0.8473\cdot\SI{200}{pF})=\SI{1.77}{k\ohm}$.
  \item Standard values (E12) inside the ranges: Standard-mode \SI{1}{k\ohm}\ldots\SI{5.6}{k\ohm}; Fast-mode
        \SI{1}{k\ohm}\ldots\SI{1.5}{k\ohm} (\SI{1.8}{k\ohm} is just above the limit, see task 4.4).
\end{steps}
\answer $R_{P,min}$ comes from the $V_{OL}$/sink current requirement (static), $R_{P,max}$ from the rise time
(dynamic). The minimum is the same for both modes, the maximum is more than three times smaller for Fast-mode.
\end{solution}

\task{S} Check the limits in the simulation: complete Table~4.3. Then increase $C_b$ to \SI{400}{pF} (the maximum
allowed by I$^2$C) with $R_P=\SI{4.7}{k\ohm}$ (model: $t_r=\SI{1.59}{\micro s}$): which mode still works?

\begin{results}{Table 4.3: low level and rise time over $R_P$ ($C_b=\SI{200}{pF}$)}{lcccccc}
  $R_P$ & \SI{5.9}{k\ohm} & \SI{4.7}{k\ohm} & \SI{1.8}{k\ohm} & \SI{1}{k\ohm} & \SI{470}{\ohm} & \SI{220}{\ohm}\\ \midrule
  $I_{low}$, model & \SI{0.56}{mA} & \SI{0.70}{mA} & \SI{1.8}{mA} & \SI{3.2}{mA} & \SI{6.6}{mA} & \SI{13.1}{mA}\\
  $V_{OL}$, model & \SI{17}{mV} & \SI{22}{mV} & \SI{56}{mV} & \SI{99}{mV} & \SI{204}{mV} & \SI{406}{mV}\\
  $V_{OL}$, simulation & \simv{17} & \simv{22} & \simv{56} & \simv{101} & \simv{211} & \simv{437}\\
  $t_r$, model & \SI{1000}{ns} & \SI{796}{ns} & \SI{305}{ns} & \SI{169}{ns} & \SI{80}{ns} & \SI{37}{ns}\\
  $t_r$, simulation & \simv{1000} & \simv{796} & \simv{306} & \simv{168} & \simv{78} & \simv{36}\\
\end{results}

\begin{solution}
\begin{steps}
  \item Model values: $I_{low}=V_{DD}/(R_P+R_{on,N})$, e.g.\ \SI{220}{\ohm}:
        $\SI{3.3}{V}/\SI{251}{\ohm}=\SI{13.1}{mA}$; $V_{OL}\approx I_{low}R_{on,N}$, e.g.\
        $\SI{13.1}{mA}\cdot\SI{30.9}{\ohm}=\SI{406}{mV}$; $t_r=0.8473\,R_PC_b$.
  \item Simulation, low level: up to \SI{1}{k\ohm} model and simulation agree; at \SI{470}{\ohm} and
        \SI{220}{\ohm} the simulation is higher (211 and \SI{437}{mV}), because $R_{on}$ grows with $V_{DS}$ in
        the triode region (the linear approximation only holds for small $V_{DS}$).
  \item Simulation, rise time: $t_r=\SI{1000}{ns}$ at \SI{5.9}{k\ohm} (exactly the Standard-mode limit) and
        $t_r=\SI{306}{ns}$ at \SI{1.8}{k\ohm}, just above the Fast-mode limit; \SI{1.5}{k\ohm} gives margin.
  \item Lower limit: our strong pin (\SI{31}{\ohm}) keeps $V_{OL}$ below \SI{0.4}{V} down to about
        \SI{250}{\ohm}. The specification only guarantees \SI{3}{mA} at \SI{0.4}{V}, so \SI{967}{\ohm} remains
        the limit for a weak but compliant device: the bus must work with every device, not only with ours.
  \item $C_b=\SI{400}{pF}$, $R_P=\SI{4.7}{k\ohm}$: $t_r=0.8473\cdot\SI{4.7}{k\ohm}\cdot\SI{400}{pF}=\SI{1.59}{\micro s}$.
        Neither mode is met; at \SI{100}{kHz} the line does not even reach $V_{DD}$ (only \SI{3.0}{V}).
  \item Required $R_P$ at \SI{400}{pF}: Standard-mode $\le\SI{1000}{ns}/(0.8473\cdot\SI{400}{pF})=\SI{2.9}{k\ohm}$
        (possible, e.g.\ \SI{2.7}{k\ohm}); Fast-mode $\le\SI{885}{\ohm}<R_{P,min}=\SI{967}{\ohm}$ (not possible
        with \SI{3}{mA} devices).
\end{steps}
\answer With \SI{400}{pF} and \SI{4.7}{k\ohm} no mode works. Standard-mode works with $R_P\le\SI{2.9}{k\ohm}$.
Fast-mode only works if all devices meet the additional rating of \SI{6}{mA} at $V_{OL}=\SI{0.6}{V}$
($R_{P,min}=\SI{450}{\ohm}$, e.g.\ \SI{820}{\ohm}), with Fast-mode Plus devices (\SI{20}{mA}), or if the bus is
split by a buffer.
\end{solution}

\task{S/D} Set the input of device~2 to \textsf{L}. What happens to the clock of device~1? Which I$^2$C feature
uses this behaviour? What would happen with two push-pull outputs (part 2)?

\begin{solution}
\begin{steps}
  \item Observation: as long as device~2 holds the line low, BUS stays low, although device~1 releases it
        (high) every half period. The clock disappears from the line.
  \item Explanation: the line is high only if all devices release it. It is the AND of all outputs
        (``wired AND''). Device~2 pulling low simply overrides the pull-up; only a small current
        $V_{DD}/R_P$ flows.
  \item Detection: device~1 can find out by reading back the line. If it released the line but reads low,
        another device is holding it.
  \item I$^2$C uses this for clock stretching (a slow slave holds SCL low until it is ready; the master waits),
        for arbitration between masters (a master that sends 1 but reads 0 has lost and stops), and for the
        acknowledge bit (the receiver pulls SDA low).
  \item Push-pull: if device~1 drove high and device~2 low, the result would be the contention of part~2:
        about \SI{28}{mA}, an undefined level of \SI{1.09}{V} and up to \SI{62}{mW} in one transistor.
\end{steps}
\answer The clock of device~1 is suppressed (wired AND). I$^2$C uses this for clock stretching, arbitration and
ACK. With push-pull outputs the same situation would be a short circuit.
\end{solution}

\task{D} While the line is low, $R_P$ dissipates $P=V_{DD}^2/R_P$ (\SI{2.3}{mW} at \SI{4.7}{k\ohm}, \SI{11}{mW} at
\SI{1}{k\ohm}). Since $t_r=0.8473\,R_PC_b$, the product $P\cdot t_r=0.8473\,V_{DD}^2C_b$ does not depend on $R_P$.
What follows for the trade-off between speed and power of an open-drain bus, compared to push-pull? Propose a
circuit that speeds up the rising edge without reducing $R_P$, and try it in the simulation (e.g.\ a current
source pull-up or a switched pull-up).

\begin{solution}
\begin{steps}
  \item Static power while low: the full supply is across $R_P$: $P=V_{DD}^2/R_P$, i.e.\
        $(\SI{3.3}{V})^2/\SI{4.7}{k\ohm}=\SI{2.3}{mW}$ and $(\SI{3.3}{V})^2/\SI{1}{k\ohm}=\SI{11}{mW}$.
  \item Product: $P\cdot t_r=\dfrac{V_{DD}^2}{R_P}\cdot0.8473\,R_PC_b=0.8473\,V_{DD}^2C_b$. $R_P$ cancels.
  \item Consequence: for a given bus ($C_b$, $V_{DD}$) speed and power can only be traded against each other.
        Halving the rise time doubles the static power.
  \item Push-pull: no static current (only $C_LV_{DD}^2f$), both edges are driven actively. It is faster and
        more efficient, but it allows no wired AND.
  \item Faster rising edge without a smaller $R_P$: (1) a constant-current pull-up charges $C_b$ linearly, so
        $t_r=0.4\,V_{DD}C_b/I$ for \SIrange{30}{70}{\percent}, and the current while low is still only $I$;
        (2) a rise-time accelerator that switches on a strong pull-up only while the line is rising above a
        threshold (e.g.\ LTC4311); (3) an active push-pull pull-up for the rising edge, as used by I$^2$C Hs-mode
        masters.
  \item In the simulation: replace $R_P$ by a current source (e.g.\ \SI{1}{mA}) and compare $t_r$ and the current
        while low.
\end{steps}
\answer An open-drain bus with a resistor can only buy speed with static power. Active pull-ups (current source,
accelerator, switched pull-up) break this trade-off, but they must never fight against a device that pulls low.
\end{solution}

%======================================================================
\needspace{0.45\textheight}
\section{Maximum output voltage: open drain with a 5\,V pull-up}\label{sec:overvoltage}
%======================================================================
\circuitlink{\linkE}{E\_open\_drain\_overvoltage.txt}

A \SI{3.3}{V} microcontroller drives a \SI{5}{V} logic input through an open-drain output and a pull-up to
\SI{5}{V}. PIN1 is a normal GPIO in open-drain mode: the PMOS is switched off, but it is still there. PIN2 is a
\SI{5}{V}-tolerant open-drain pin without a PMOS. The \SI{3.3}{V} regulator can only source current (diode).
The resistor ``MCU load'' represents the current consumption of the chip.

\schematic[0.7]{E}

\task{P/S} Table~5.1 gives the expected high levels and the current that flows from the \SI{5}{V} supply into the
\SI{3.3}{V} supply via PIN1 ($R_{PU}=\SI{4.7}{k\ohm}$, diode forward voltage about \SI{0.7}{V}). Explain the
values, then measure them.

\begin{results}{Table 5.1: open-drain pins with a pull-up to \SI{5}{V} ($R_{PU}=\SI{4.7}{k\ohm}$)}{lccc}
  & high level PIN1 & high level PIN2 & injection current\\ \midrule
  model & $V_{DD}+V_F\approx\SI{4.0}{V}$ & \SI{5.0}{V} & $\dfrac{5-4.0}{4.7}\,\si{mA}\approx\SI{0.2}{mA}$\\
  simulation & \simv{4.10} & \simv{5.00} & \simv{0.19}\\
\end{results}

\begin{solution}
\begin{steps}
  \item PIN2 has no PMOS and therefore no diode to $V_{DD}$. When the NMOS is off, no current flows through the
        pull-up and the pin rises to the full \SI{5.0}{V}.
  \item PIN1: the PMOS is switched off, but its body diode (pin $\rightarrow$ $V_{DD}$) is still there. As soon as
        the pin rises above $V_{DD}+V_F$, the diode conducts and clamps the pin to
        $\SI{3.3}{V}+\SI{0.7}{V}\approx\SI{4.0}{V}$.
  \item Injection current: the rest of the \SI{5}{V} drops across the pull-up:
        $I=(\SI{5}{V}-\SI{4.0}{V})/\SI{4.7}{k\ohm}\approx\SI{0.2}{mA}$. It flows through the body diode into
        the \SI{3.3}{V} supply.
  \item Simulation: PIN1 \SI{4.10}{V}, PIN2 \SI{5.00}{V},
        injection current $(5-4.10)\,\si{V}/\SI{4.7}{k\ohm}=\SI{0.19}{mA}$. The clamp level is slightly
        higher than the estimate with \SI{0.7}{V}.
\end{steps}
\answer The normal GPIO is clamped to about one diode drop above $V_{DD}$ and feeds \SI{0.2}{mA} into the
\SI{3.3}{V} rail; the \SI{5}{V}-tolerant pin reaches \SI{5}{V} without any current.
\end{solution}

\task{D} Is the high level of PIN1 sufficient for a \SI{5}{V} CMOS input with $V_{IH}=0.7\cdot\SI{5}{V}$? Is it
allowed?

\begin{solution}
\begin{steps}
  \item Required high level: $V_{IH}=0.7\cdot\SI{5}{V}=\SI{3.5}{V}$.
  \item PIN1 reaches \SI{4.10}{V} $>$ \SI{3.5}{V}: logically sufficient, with \SI{0.6}{V} margin.
  \item Absolute maximum rating: $IOVDD+\SI{0.5}{V}=\SI{3.8}{V}$. PIN1 is \SI{0.3}{V} above it, and a current
        flows through the body diode.
  \item The margin is not reliable: the diode voltage drops with temperature (about \SI{-2}{mV/K}) and with a
        smaller current (larger $R_{PU}$), and $V_{DD}$ itself has a tolerance.
\end{steps}
\answer Logically yes, but it is not allowed: the pin is operated outside its absolute maximum rating, with a
permanent injection current.
\end{solution}

\task{S/D} Set $R_{PU}=\SI{1}{k\ohm}$ and keep the MCU load at \SI{33}{k\ohm} (\SI{100}{\micro A}, e.g.\ in a
sleep mode). Observe $V_{DD}$ (simulation: \underline{\hspace{2cm}}\simv{3.9\,V}) and explain the effect
(\emph{back-powering}): the regulator can only source current, so $V_{DD}$ rises until the chip consumes the
injected current. What happens if the microcontroller is switched off ($V_{DD}=0$) but the \SI{5}{V} side stays on?

\begin{solution}
\begin{steps}
  \item Observation: $V_{DD}$ rises to \SI{3.9}{V} (with \SI{470}{\ohm}: \SI{4.0}{V}), above the absolute
        maximum supply voltage of the chip.
  \item Current balance at the \SI{3.3}{V} rail: current flows in from the regulator (only in one direction,
        diode) and through the body diode of PIN1, and out through the MCU load.
  \item Injected current: while PIN1 is high, \SI{0.25}{mA} flow through the body diode into the rail; since
        PIN1 is high half of the time, the average is \SI{0.12}{mA}.
  \item Equilibrium: the load draws only $V_{DD}/\SI{33}{k\ohm}$. The regulator cannot sink the excess, so
        $V_{DD}$ rises until the load consumes the injected current:
        $\SI{3.9}{V}/\SI{33}{k\ohm}=\SI{0.12}{mA}$. The regulator then delivers nothing.
  \item Cross-check: with a load of \SI{1}{mA} (\SI{3.3}{k\ohm}) the chip consumes more than is injected; the
        regulator delivers the rest and the rail stays at \SI{3.27}{V}. Back-powering is a problem especially in
        sleep modes.
  \item Microcontroller off ($V_{DD}=0$): the chip is now powered only through its pin (``phantom powering'').
        It may start partially, hold the bus low, draw a large injection current, or come up in an undefined
        state when its supply returns.
\end{steps}
\answer The injected current back-powers the \SI{3.3}{V} rail when the chip consumes less than is injected. A
switched-off microcontroller is powered through its pin. Data sheets therefore forbid voltages above
$V_{DD}+\SI{0.3}{V}\ldots\SI{0.5}{V}$ on non-tolerant pins, also when the chip is off.
\end{solution}

\task{D} Look up in the data sheets of an RP2040, an RP2350 and an STM32F4 which pins are \SI{5}{V} tolerant
and how large the injection current per pin and in total may be. Name two circuits that shift a
\SI{3.3}{V} open-drain signal to \SI{5}{V} correctly, one of them bidirectional (for I$^2$C).

\begin{solution}
\begin{steps}
  \item RP2040: no GPIO is \SI{5}{V} tolerant; the pin voltage is limited to $IOVDD+\SI{0.5}{V}$ (absolute
        maximum).
  \item RP2350: the GPIOs are \SI{5}{V} tolerant while IOVDD is powered (and fail-safe when unpowered); the data
        sheet lists the exact conditions (e.g.\ not for pins used as ADC inputs).
  \item STM32F4: pins marked ``FT'' in the pin table are \SI{5}{V} tolerant (in input or open-drain mode). The
        injection current is limited to about $\pm\SI{5}{mA}$ per pin and $\pm\SI{25}{mA}$ in total; positive
        injection is not allowed on FT pins.
  \item Level shifter 1 (unidirectional): a \SI{5}{V}-tolerant pin in open-drain mode with a pull-up to
        \SI{5}{V}, or an external open-drain transistor (part~\ref{sec:oc}).
  \item Level shifter 2 (bidirectional, I$^2$C): one N-channel MOSFET (e.g.\ BSS138) per line, gate at
        \SI{3.3}{V}, source on the \SI{3.3}{V} side, drain on the \SI{5}{V} side, pull-ups on both sides. If the
        \SI{3.3}{V} side pulls low, the MOSFET turns on and pulls the \SI{5}{V} side low; if the \SI{5}{V} side
        pulls low, the body diode and then the channel pull the \SI{3.3}{V} side low. Alternatively dedicated
        translators (e.g.\ PCA9306, TXS0102).
\end{steps}
\answer RP2040: no \SI{5}{V} tolerant pins; RP2350: tolerant while powered; STM32F4: FT pins. A BSS138 per line
with pull-ups on both sides shifts I$^2$C correctly in both directions.
\end{solution}

%======================================================================
\needspace{0.45\textheight}
\section{Open collector: high voltage, high current, inductive loads}\label{sec:oc}
%======================================================================
\circuitlink{\linkF}{F\_open\_collector\_inductive.txt}

A \SI{3.3}{V} GPIO switches an NPN transistor ($\beta=100$) through the base resistor $R_B$. The collector is an
open-collector output for a \SI{12}{V} load: a relay coil (\SI{0.1}{H}, \SI{400}{\ohm}, \SI{1}{nF} winding
capacitance) or a \SI{12}{V}/\SI{5}{W} lamp. The switch \emph{flyback} adds a free-wheeling diode, opening the
switch \emph{bypass} puts a \SI{24}{V} Zener diode in series with it. The timescale is shortened (switching at
\SI{200}{Hz}) so that you can see several cycles.

\schematic[0.62]{F}

\task{P} Table~6.1 gives the operating point with $R_B=\SI{1}{k\ohm}$. Is the transistor saturated?

\begin{results}{Table 6.1: operating point, relay coil, $R_B=\SI{1}{k\ohm}$}{lccccc}
  & $I_C=\dfrac{\SI{12}{V}}{\SI{400}{\ohm}}$ & $I_B=\dfrac{3.3-0.7}{R_B}$ & $\beta I_B$ & overdrive $\dfrac{\beta I_B}{I_C}$ & $V_{CE}$\\ \midrule
  model & \SI{30}{mA} & \SI{2.6}{mA} & \SI{260}{mA} & 8.7 & $<\SI{0.2}{V}$\\
  simulation & \simv{29.8} & \simv{2.61} & & & \simv{0.07}\\
\end{results}

\begin{solution}
\begin{steps}
  \item Required collector current: if the transistor were an ideal switch, the coil would see the full
        \SI{12}{V}: $I_C=\SI{12}{V}/\SI{400}{\ohm}=\SI{30}{mA}$.
  \item Base current: the GPIO delivers \SI{3.3}{V}, the base-emitter diode takes \SI{0.7}{V}:
        $I_B=(3.3-0.7)\,\si{V}/\SI{1}{k\ohm}=\SI{2.6}{mA}$ (the GPIO can deliver that easily, part~1).
  \item Possible collector current in the active region: $\beta I_B=100\cdot\SI{2.6}{mA}=\SI{260}{mA}$.
  \item Overdrive: $\beta I_B/I_C=260/30=8.7>1$. The transistor could conduct far more than the load allows, so
        $V_{CE}$ drops to the saturation voltage (simulation: \SI{0.07}{V}).
\end{steps}
\answer Yes, the transistor is saturated with an overdrive factor of 8.7; the simulation shows $I_C=\SI{29.8}{mA}$
and $V_{CE}=\SI{0.07}{V}$.
\end{solution}

\task{S/D} Without flyback diode: when the transistor switches off, the coil current continues and charges the
winding capacitance. Energy conservation ($\tfrac12LI^2=\tfrac12CV^2$) gives a peak of
$\SI{12}{V}+I\sqrt{L/C}=\SI{12}{V}+\SI{300}{V}$. Measure it (simulation: \underline{\hspace{2cm}}\simv{296\,V}) and compare with
$V_{CEO}$ of a BC547 (\SI{45}{V}) or of a ULN2003 driver (\SI{50}{V}). What does the simulation not show?

\begin{solution}
\begin{steps}
  \item Before switching off, the coil stores $E=\tfrac12LI^2=\tfrac12\cdot\SI{0.1}{H}\cdot(\SI{30}{mA})^2=\SI{45}{\micro J}$.
  \item The coil current cannot change abruptly. With the transistor off its only path is the winding
        capacitance, which is charged until all energy has moved into it: $\tfrac12LI^2=\tfrac12C\,\Delta V^2$.
  \item $\Delta V=I\sqrt{L/C}=\SI{30}{mA}\cdot\sqrt{\SI{0.1}{H}/\SI{1}{nF}}=\SI{30}{mA}\cdot\SI{10}{k\ohm}=\SI{300}{V}$
        above the \SI{12}{V} supply, i.e.\ about \SI{312}{V} at the collector.
  \item Simulation: peak $V_C=\SI{296}{V}$ (a little energy is lost in the coil resistance), followed by a
        ringing that drives the collector to \SI{-14}{V}, where the base-collector diode conducts.
  \item Comparison: $V_{CEO}=\SI{45}{V}$ (BC547) or \SI{50}{V} (ULN2003) is exceeded six times.
  \item What the simulation does not show: the transistor model has no breakdown. A real transistor breaks down
        (avalanche) at about \SIrange{45}{80}{V} and absorbs the coil energy in a small volume, which degrades or
        destroys it. Neither the heating nor the EMI of the fast edge, nor the coupling back into the GPIO via
        the base, is modelled; the real winding capacitance and damping determine the actual peak.
\end{steps}
\answer Without a flyback diode the collector voltage reaches about \SI{300}{V}, far above $V_{CEO}$. In reality
the transistor would break down and be damaged; the simulation only shows the voltage the circuit would generate.
\end{solution}

\task{S} Close \emph{flyback} and measure the maximum collector voltage and the time until the coil current has
fallen to \SI{10}{\percent} (right-click the scope, \emph{Properties}, to change the time base). Then open
\emph{bypass} (diode + \SI{24}{V} Zener) and measure again. Explain why the Zener variant is faster. Why does that
matter for a relay?

\begin{results}{Table 6.2: switching off the relay coil ($\tau=L/R=\SI{0.25}{ms}$)}{lcccc}
  & \multicolumn{2}{c}{diode} & \multicolumn{2}{c}{diode + \SI{24}{V} Zener}\\
  & $V_{C,max}$ & $t_{10\%}$ & $V_{C,max}$ & $t_{10\%}$\\ \midrule
  model & $12+0.7=\SI{12.7}{V}$ & \SI{475}{\micro s} & $12+24+0.7=\SI{36.7}{V}$ & \SI{87}{\micro s}\\
  simulation & \simv{12.8} & \simv{474} & \simv{36.8} & \simv{87}\\
\end{results}
{\small Model: $i(t)=-V_{cl}/R+(I_0+V_{cl}/R)\,e^{-t/\tau}$ with the clamp voltage $V_{cl}=V_F$ or $V_Z+V_F$.}

\begin{solution}
\begin{steps}
  \item Diode only, maximum voltage: the coil current continues through the diode to the \SI{12}{V} rail, so the
        collector rises only one diode drop above it: $V_{C,max}=\SI{12.7}{V}$ (simulation \SI{12.8}{V}).
  \item Decay: in the loop coil--diode, $L\,di/dt=-Ri-V_{cl}$ with $V_{cl}=V_F$. The solution is
        $i(t)=-V_{cl}/R+(I_0+V_{cl}/R)\,e^{-t/\tau}$ with $\tau=L/R=\SI{0.1}{H}/\SI{400}{\ohm}=\SI{0.25}{ms}$.
  \item Time to \SI{10}{\percent} (\SI{3}{mA}) with $V_{cl}/R=\SI{0.7}{V}/\SI{400}{\ohm}=\SI{1.75}{mA}$:
        $\SI{3}{mA}=-\SI{1.75}{mA}+\SI{31.75}{mA}\,e^{-t/\tau}\Rightarrow t=\tau\ln(31.75/4.75)=\SI{475}{\micro s}$
        (simulation \SI{474}{\micro s}; without $V_F$ it would be $\tau\ln10=\SI{576}{\micro s}$).
  \item Diode + Zener, maximum voltage: the current now has to flow against $V_Z+V_F$:
        $V_{C,max}=12+24+0.7=\SI{36.7}{V}$ (simulation \SI{36.8}{V}), still below $V_{CEO}=\SI{45}{V}$.
  \item Decay with $V_{cl}=\SI{24.7}{V}$, $V_{cl}/R=\SI{61.8}{mA}$:
        $i(t)=-\SI{61.8}{mA}+\SI{91.8}{mA}\,e^{-t/\tau}=\SI{3}{mA}\Rightarrow t=\tau\ln(91.8/64.8)=\SI{87}{\micro s}$
        (simulation \SI{87}{\micro s}).
  \item Why faster: the energy $\tfrac12LI^2$ is removed with the power $P=V_{cl}\cdot i$. With \SI{24.7}{V}
        instead of \SI{0.7}{V} across the clamp, the coil voltage $L\,di/dt$ is much larger, so the current falls
        much faster.
\end{steps}
\answer The Zener variant releases the relay about five times faster (\SI{87}{\micro s} instead of
\SI{475}{\micro s}) at the cost of a higher collector voltage. For a relay the current decay determines how fast
the contacts open: a slow decay lets the contacts drift apart slowly, which increases arcing and contact wear. The
designer trades transistor voltage rating against release speed.
\end{solution}

\task{S/D} Increase $R_B$ to \SI{47}{k\ohm}: $I_B=\SI{55}{\micro A}$, so $I_C\le\beta I_B=\SI{5.5}{mA}$ and
$V_{CE}\approx 12-5.5\cdot0.4=\SI{9.8}{V}$. Measure $I_C$ and $V_{CE}$
(simulation: \underline{\hspace{1.5cm}}\simv{5.7\,mA} and \underline{\hspace{1.5cm}}\simv{9.7\,V}). Will the relay still pull in?

\begin{solution}
\begin{steps}
  \item Base current: $I_B=(3.3-0.7)\,\si{V}/\SI{47}{k\ohm}=\SI{55}{\micro A}$.
  \item Maximum collector current: $\beta I_B=\SI{5.5}{mA}<\SI{30}{mA}$ required: overdrive $0.18<1$, so the
        transistor is not saturated but in the active region (current source).
  \item Collector voltage: the coil drops $\SI{5.5}{mA}\cdot\SI{400}{\ohm}=\SI{2.2}{V}$, so
        $V_{CE}=\SI{12}{V}-\SI{2.2}{V}=\SI{9.8}{V}$.
  \item Simulation: $I_C=\SI{5.7}{mA}$, $V_{CE}=\SI{9.7}{V}$. The transistor dissipates
        $\SI{9.7}{V}\cdot\SI{5.7}{mA}=\SI{55}{mW}$ instead of about \SI{2}{mW} when saturated.
  \item The coil receives $5.7/30=\SI{19}{\percent}$ of its nominal current; a relay typically needs
        \SIrange{70}{75}{\percent} to pull in.
\end{steps}
\answer No, the relay will not pull in. Rule: choose $R_B$ so that $\beta_{min}I_B\ge 2\ldots5\cdot I_C$
(overdrive), using the minimum $\beta$ of the data sheet, not the typical one.
\end{solution}

\task{S/D} Switch the load to the lamp (\SI{12}{V}/\SI{5}{W}, needs \SI{414}{mA}) and complete Table~6.3. A
small-signal transistor in SOT-23 tolerates about \SI{250}{mW}. Why is \SI{330}{\ohm} still not a good solution?
Propose a better driver.

\begin{results}{Table 6.3: lamp load}{lcccc}
  & $I_B$ & $I_C$ & $V_{CE}$ & $P=V_{CE}I_C$\\ \midrule
  $R_B=\SI{1}{k\ohm}$, model & \SI{2.6}{mA} & $\beta I_B=\SI{260}{mA}$ & $12-0.26\cdot29=\SI{4.5}{V}$ & \SI{1.2}{W}\\
  $R_B=\SI{1}{k\ohm}$, simulation & \simv{2.56} & \simv{256} & \simv{4.57} & \simv{1.17}\\
  $R_B=\SI{330}{\ohm}$, model & \SI{7.9}{mA} & \SI{414}{mA} (saturated) & $\approx\SI{0.1}{V}$ & $\approx\SI{50}{mW}$\\
  $R_B=\SI{330}{\ohm}$, simulation & \simv{7.72} & \simv{410} & \simv{0.12} & \simv{0.05}\\
\end{results}

\begin{solution}
\begin{steps}
  \item Lamp: $I=P/V=\SI{5}{W}/\SI{12}{V}\approx\SI{414}{mA}$, hot resistance $R=V^2/P=\SI{28.8}{\ohm}$.
  \item $R_B=\SI{1}{k\ohm}$: $I_B=\SI{2.6}{mA}$, $\beta I_B=\SI{260}{mA}<\SI{414}{mA}$: active region. The lamp
        drops $\SI{0.26}{A}\cdot\SI{28.8}{\ohm}\approx\SI{7.5}{V}$, so $V_{CE}=\SI{4.5}{V}$ and
        $P=\SI{4.5}{V}\cdot\SI{260}{mA}=\SI{1.2}{W}$. Simulation: \SI{256}{mA}, \SI{4.57}{V}, \SI{1.17}{W}.
        The lamp stays dim, and the transistor (\SI{250}{mW}) overheats within seconds.
  \item $R_B=\SI{330}{\ohm}$: $I_B=\SI{2.6}{V}/\SI{330}{\ohm}=\SI{7.9}{mA}$, $\beta I_B=\SI{790}{mA}>\SI{414}{mA}$:
        saturated, $V_{CE}\approx\SI{0.1}{V}$, $P\approx\SI{0.12}{V}\cdot\SI{0.41}{A}=\SI{50}{mW}$. Simulation:
        \SI{7.72}{mA}, \SI{410}{mA}, \SI{0.12}{V}, \SI{0.05}{W}.
  \item Why \SI{330}{\ohm} is still not good: (1) \SI{7.7}{mA} is a substantial load for the GPIO; $V_{OH}$ drops
        (part~1) and the per-domain budget shrinks. (2) The overdrive is only $790/414=1.9$ with the typical
        $\beta$; with the minimum $\beta$ of the data sheet, and since $\beta$ falls at high currents, the
        transistor may leave saturation. (3) A cold filament draws about ten times the nominal current when
        switched on.
\end{steps}
\answer A bipolar small-signal transistor is the wrong component for this load. Better: a logic-level N-MOSFET
(e.g.\ with $R_{DS(on)}<\SI{0.1}{\ohm}$ at $V_{GS}=\SI{2.5}{V}$, gate resistor about \SI{100}{\ohm}, pull-down
\SI{100}{k\ohm}), which needs no static gate current; a Darlington array such as the ULN2003 (\SI{500}{mA},
integrated clamp diodes); or a high-side switch with current limit.
\end{solution}

\task{D} Summarise: what limits the maximum output voltage and the maximum output current of a push-pull pin, of
an open-drain pin and of an external open-collector stage?

\begin{solution}
\begin{steps}
  \item Push-pull pin: the output voltage cannot exceed $V_{DD}$ (the PMOS connects to $V_{DD}$); an external
        voltage is clamped at $V_{DD}+V_F$ by the body diodes (part~1, part~\ref{sec:overvoltage}). The current is
        limited by $R_{on}$ (level drop), by $I_{sat}$ (short circuit), and by the per-pin and per-domain ratings
        (heat, electromigration).
  \item Open-drain pin: the high voltage is set by the pull-up, but limited by the pin structure: body/ESD diode
        to $V_{DD}$ on a normal pin, about \SI{5.5}{V} on a \SI{5}{V}-tolerant pin. The sink current is limited
        by the NMOS as for push-pull; the high level is only as strong as the pull-up (slow edges, part~4).
  \item External open collector/drain: the voltage is limited by $V_{CEO}$ or $V_{DS,max}$ of the transistor,
        including inductive kick-back (part~\ref{sec:oc}). The current is limited by $\beta I_B$ (or $V_{GS}$ and
        $R_{DS(on)}$ for a MOSFET), the saturation voltage and the power dissipation of the package.
\end{steps}
\answer Push-pull: $V_{DD}$ and the on-chip transistors; open drain: the pin structure and the NMOS; external
stage: the ratings of the external transistor, which can be chosen for the load.
\end{solution}

%======================================================================
\needspace{0.45\textheight}
\section{Fast edges on a PCB trace}
%======================================================================
\circuitlink{\linkG}{G\_transmission\_line.txt}

The push-pull pin drives a \SI{50}{\ohm} trace with \SI{1}{ns} delay (about \SI{15}{cm} on FR4). At the end there is
the input of another chip: \SI{5}{pF} and ESD diodes to $V_{DD}$ and GND. $R_S$ is a series resistor at the
driver. The output stage is set to its highest drive strength and fastest slew rate, modelled by
$\beta_N=\SI{30}{mA/V^2}$ and $\beta_P=\SI{25}{mA/V^2}$. The slew rate is adjustable as on many
microcontrollers (RP2040: \texttt{SLEWFAST}, STM32: \texttt{OSPEEDR}).

\schematic[0.75]{G}

\task{P} A trace is ``electrically long'' if the edge is shorter than about twice its delay. Table~7.1 gives the
values for $R_S=0$ and the fastest slew rate. Explain the sequence: launched step, doubling at the open end,
clamping, reflection at the driver.

\begin{results}{Table 7.1: reflections for $R_S=0$ ($R_{on,P}=1/(\SI{25}{mA/V^2}\cdot\SI{2.7}{V})=\SI{15}{\ohm}$)}{lc}
  edge time / round trip & $<\SI{1}{ns}$ / \SI{2}{ns}: electrically long\\
  $\Gamma_S=(R_{on}+R_S-Z_0)/(R_{on}+R_S+Z_0)$ & $(15-50)/(15+50)=-0.54$\\
  $\Gamma_L$ (open end, \SI{5}{pF} only rounds the edge) & $\approx+1$\\
  launched step $V_{DD}\,Z_0/(Z_0+R_{on})$ & \SI{2.54}{V}\\
  first voltage at RX (before clamping) & $2\cdot\SI{2.54}{V}=\SI{5.1}{V}$\\
  matched series resistor $R_S=Z_0-R_{on}$ & $\approx\SI{36}{\ohm}$ (\SI{33}{\ohm} standard value)\\
\end{results}

\begin{solution}
\begin{steps}
  \item Electrically long: in its fastest setting the output edge takes well below \SI{1}{ns}; the round trip
        (there and back) takes $2\cdot\SI{1}{ns}=\SI{2}{ns}$. The edge is shorter than the round trip, so the
        line voltage is not the same everywhere: the trace must be treated as a transmission line.
  \item Source impedance: the strong PMOS has $R_{on,P}=1/(\SI{25}{mA/V^2}\cdot\SI{2.7}{V})=\SI{15}{\ohm}$
        (NMOS: \SI{12}{\ohm}).
  \item Launched step: at the first moment the driver only sees $Z_0$, a voltage divider of $R_{on}$ and $Z_0$:
        $\SI{3.3}{V}\cdot50/(50+15)=\SI{2.54}{V}$ travels down the line.
  \item Open end: the input (\SI{5}{pF}) is nearly open, $\Gamma_L\approx+1$. The wave is reflected with the same
        sign and the voltage at RX doubles to $2\cdot\SI{2.54}{V}=\SI{5.1}{V}$.
  \item Clamping: \SI{5.1}{V} is \SI{1.8}{V} above $V_{DD}$, so the ESD diode to $V_{DD}$ conducts and limits RX
        to about $V_{DD}+\SI{0.6}{V}$. The diode current flows into the supply of the receiver.
  \item Reflection at the driver: the returning wave sees $\Gamma_S=(15-50)/(15+50)=-0.54$. It is reflected with
        opposite sign and arrives at RX \SI{2}{ns} later as an undershoot of roughly $-0.54\cdot\SI{2.5}{V}$, which
        is again limited by the ESD diode to GND.
\end{steps}
\answer The edge rings with a period of \SI{4}{ns} (two round trips) and decays by the factor $\Gamma_S$ per
round trip, while the ESD diodes clamp the peaks.
\end{solution}

\task{S} Measure overshoot and undershoot at RX and the peak currents in the ESD diodes for the values of $R_S$ in
Table~7.2. Which $R_S$ gives the cleanest edge? Compare with $R_S=Z_0-R_{on}$.

\begin{results}{Table 7.2: series termination (fastest slew rate)}{lccccc}
  $R_S$ & \SI{0}{\ohm} & \SI{10}{\ohm} & \SI{22}{\ohm} & \SI{36}{\ohm} & \SI{50}{\ohm}\\ \midrule
  $V_{RX,max}$, simulation & \simv{3.92} & \simv{3.89} & \simv{3.74} & \simv{3.30} & \simv{3.30}\\
  $V_{RX,min}$, simulation & \simv{-0.64} & \simv{-0.61} & \simv{-0.52} & \simv{-0.02} & \simv{0.00}\\
  $I_{ESD}$ to $V_{DD}$ / GND, simulation & \simv{30/40\,mA} & \simv{14/21} & \simv{1/4} & \simv{0} & \simv{0}\\
\end{results}

\begin{solution}
\begin{steps}
  \item $R_S=0$: RX swings to \SI{3.9}{V} and \SI{-0.64}{V}, beyond the \SI{\pm0.5}{V} limits. The ESD diodes
        conduct up to \SI{30}{mA} (to $V_{DD}$) and \SI{40}{mA} (to GND) on every edge, and the ringing can cross
        the input thresholds several times (double clocking).
  \item Measured values for all $R_S$:\\[1mm]
\begin{tabular}{lccccc}
  \toprule
  $R_S$ / \si{\ohm} & 0 & 10 & 22 & 36 & 50\\
  \midrule
  $V_{RX,max}$ / V & 3.92 & 3.89 & 3.74 & 3.30 & 3.30\\
  $V_{RX,min}$ / V & $-0.64$ & $-0.61$ & $-0.52$ & $-0.02$ & 0.00\\
  $I_{ESD}$ / mA & 30/40 & 14/21 & 1/4 & 0 & 0\\
  \bottomrule
\end{tabular}
  \item With increasing $R_S$ the launched step becomes smaller and $\Gamma_S$ approaches zero, so overshoot and
        ESD currents decrease. At \SI{36}{\ohm} they vanish.
  \item Theory: with $R_S+R_{on}=Z_0$ the driver launches $V_{DD}\cdot Z_0/(2Z_0)=V_{DD}/2$; the open end doubles
        it to exactly $V_{DD}$; the reflected wave arrives at a matched source ($\Gamma_S=0$) and is absorbed.
        $R_S=Z_0-R_{on}$ lies between $50-15=\SI{35}{\ohm}$ (rising) and $50-12=\SI{38}{\ohm}$ (falling edge).
  \item $R_S=\SI{50}{\ohm}$ is also clean, but the source is now overdamped: the step at RX reaches $V_{DD}$ only
        in several round trips, the edge becomes slower.
\end{steps}
\answer The cleanest edge is at $R_S\approx\SI{36}{\ohm}=Z_0-R_{on}$ (standard value \SI{33}{\ohm}), exactly as
predicted by series termination. Larger values only slow the edge.
\end{solution}

\task{S/D} Set $R_S$ back to 0 and reduce the slew rate. From which slew rate on does the ringing disappear?
Why should you select the slowest slew rate that meets the timing of your interface?

\begin{solution}
\begin{steps}
  \item The slider sets the slew rate of the gate; the output switches during only part of the gate swing, so
        measure the edge at PIN, not at the gate.
  \item \SI{1}{V/ns}: the ringing is almost unchanged (ESD currents \SI{19}{mA}/\SI{24}{mA}).
  \item \SI{0.3}{V/ns}: the ringing is reduced (\SI{4.5}{mA}/\SI{6.6}{mA}).
  \item \SI{0.1}{V/ns}: the output edge takes several \si{ns}, much longer than the \SI{2}{ns} round trip. The
        reflections arrive while the edge is still rising and merge with it; the trace behaves like a lumped
        capacitance and the edge is clean.
  \item Criterion: the ringing disappears when the edge time is several times the round-trip time
        ($t_r\gtrsim 2\ldots3\cdot 2t_d$).
\end{steps}
\answer The ringing disappears at about \SI{0.1}{V/ns}. Choose the slowest slew rate that meets the timing: a
slow edge has a narrower spectrum, which means less ringing, less crosstalk, less radiated EMI, smaller ground
bounce when many pins switch together, and less supply current (part~3).
\end{solution}

%======================================================================
\needspace{0.45\textheight}
\section{Design task}
%======================================================================

\task{P/S/D} A \SI{3.3}{V} microcontroller board has to
\begin{enumerate}[label=(\alph*)]
  \item read a \SI{5}{V} sensor over I$^2$C at \SI{400}{kHz}; the cable is \SI{30}{cm} long, the total bus
        capacitance is estimated at \SI{250}{pF},
  \item switch a \SI{24}{V}/\SI{50}{mA} solenoid valve at up to \SI{10}{Hz},
  \item send a \SI{10}{MHz} SPI clock over a \SI{10}{cm} trace to a flash memory.
\end{enumerate}
For each signal choose the output type and configuration of the pin (push-pull/open drain, drive strength, slew
rate), dimension the external components and verify your choice in a simulation that you build from the given
circuits. Document the remaining risks.

\begin{solution}
A possible solution (other justified solutions are equally valid):
\begin{enumerate}[label=(\alph*)]
  \item I$^2$C sensor at \SI{5}{V}, \SI{400}{kHz}, \SI{250}{pF}:
        \begin{steps}
          \item Output type: open drain is mandatory (wired AND, clock stretching, ACK).
          \item Level: the RP2040 pins are not \SI{5}{V} tolerant $\Rightarrow$ a BSS138 level shifter per line
                (or a PCA9306) with pull-ups on both sides (task 5.4).
          \item Pull-up on the \SI{5}{V} side, where the cable (\SI{250}{pF}) is:
                $R_{P,max}=\SI{300}{ns}/(0.8473\cdot\SI{250}{pF})=\SI{1.42}{k\ohm}$, but
                $R_{P,min}=(\SI{5}{V}-\SI{0.4}{V})/\SI{3}{mA}=\SI{1.53}{k\ohm}$: no resistor meets both.
          \item Options: devices rated for \SI{6}{mA} at \SI{0.6}{V} ($R_{P,min}=\SI{4.4}{V}/\SI{6}{mA}=\SI{733}{\ohm}$,
                choose \SI{1.2}{k\ohm}: $t_r=\SI{254}{ns}$); a rise-time accelerator or bus buffer; or
                Standard-mode at \SI{100}{kHz} ($R_{P,max}=\SI{4.7}{k\ohm}$, choose
                \SI{1.8}{k\ohm}\ldots\SI{4.7}{k\ohm}). On the short \SI{3.3}{V} side \SI{4.7}{k\ohm} is enough.
          \item Simulation: part~4 with a \SI{5}{V} pull-up and \SI{250}{pF}; check $t_r$ and $V_{OL}$.
          \item Risk: the capacitance is only an estimate; measure $t_r$ on the real board.
        \end{steps}
  \item Solenoid valve \SI{24}{V}/\SI{50}{mA} at \SI{10}{Hz}:
        \begin{steps}
          \item Coil resistance $R=\SI{24}{V}/\SI{50}{mA}=\SI{480}{\ohm}$. The GPIO cannot switch \SI{24}{V}
                $\Rightarrow$ external low-side switch: logic-level N-MOSFET or a ULN2003 channel.
          \item Pin: push-pull, any drive strength, slow slew rate (\SI{10}{Hz} needs no fast edges); gate
                resistor \SI{100}{\ohm}, pull-down \SI{100}{k\ohm} so that the valve stays off during reset (the
                pins are inputs then).
          \item Free-wheeling: diode + \SI{24}{V} Zener if the valve must close quickly; then
                $V_{DS,max}\ge24+24+1\approx\SI{50}{V}$, choose a \SI{60}{V} MOSFET. Otherwise a diode alone
                ($V_{DS}\ge\SI{40}{V}$ with margin).
          \item Simulation: part~6 with \SI{24}{V}, \SI{480}{\ohm} and the valve inductance.
          \item Risk: the valve inductance is unknown without a data sheet; measure the kick-back on the board.
        \end{steps}
  \item SPI clock \SI{10}{MHz} over \SI{10}{cm}:
        \begin{steps}
          \item Output type: push-pull (open drain would be far too slow, part~4).
          \item Trace: \SI{10}{cm} on FR4 has about \SI{0.67}{ns} delay, \SI{1.3}{ns} round trip. With the fastest
                edge it is electrically long (part~7).
          \item Edge budget: the half period is \SI{50}{ns}, so an edge of a few ns is entirely sufficient.
          \item Configuration: moderate drive strength and slow slew rate, plus a series resistor
                $R_S=Z_0-R_{on}$ (e.g.\ \SI{33}{\ohm}) close to the MCU.
          \item Simulation: part~7 with the trace delay changed to \SI{0.67}{ns}.
          \item Risk: ringing at the flash causes double clocking and corrupted data.
        \end{steps}
\end{enumerate}
\end{solution}

\end{document}
